%%%%%%%%%%%%%%%%%%%%%%%%%%%%%%%%%%%%%%%%%%%%%%%%%%%%%%%%%%%%%%%%%%%%%%%%%%%%%%

%\subsection{A polynomial bound on mixing sequences}
%\label{sec: polynomial - arbitrary n}
%%%%%%%%%%%%%%%%%%%%%%%%%%%%%%%%%%%%%%%%%%%%%%%%%%%%%%%%%%%%%%%%%%%%%%%%%%%%%%%

%\subsection{A polynomial bound on mixing sequences}
%\label{sec: polynomial - arbitrary n}
%%%%%%%%%%%%%%%%%%%%%%%%%%%%%%%%%%%%%%%%%%%%%%%%%%%%%%%%%%%%%%%%%%%%%%%%%%%%%%%

%\subsection{A polynomial bound on mixing sequences}
%\label{sec: polynomial - arbitrary n}
%%%%%%%%%%%%%%%%%%%%%%%%%%%%%%%%%%%%%%%%%%%%%%%%%%%%%%%%%%%%%%%%%%%%%%%%%%%%%%%

%\subsection{A polynomial bound on mixing sequences}
%\label{sec: polynomial - arbitrary n}
%\input{6-4_polynomial_bound.tex}

We now complete the proof of Theorem~\ref{thm: miscibility characterization}(b).
In Section~\ref{sec: sufficiency of condition mc} we have already established an existence of a perfect-mixing 
sequence for $C$ with intermediate precision $1$.
It still remains to show that this mixing sequence can be modified to have length that is polynomial in $s(C)$.

As explained in the beginning of Section~\ref{sec: polynomial proof}, it remains to
give such a bound for the configuration $E$ constructed from $C$ in Section~\ref{subsec: proof for n greater-equal than 7}
and for mixing sequences with precision $0$ (that is, with only integral concentrations).
Recall that in Section~\ref{sec: polynomial - near final} we gave a polynomial bound if $E$ is near-final.
In Section~\ref{sec: an exponential bound} we gave a bound for arbitrary configurations $E$ that is
exponential only in $n$, thus yielding a polynomial bound for constant $n$.
We will be using these properties in our construction in this section. 
If $E$ is not near-final and $n$ is arbitrary, the general idea for achieving a polynomial bound
is similar to that from Sections~\ref{sec: polynomial - near final} and~\ref{sec: an exponential bound}:
we try to mix a pair of concentrations that are far apart,
to guarantee that we make a quick progress, as measured by the decrease of $\potential(E)$. 
(This is captured by Observation~\ref{obs: far-appart mix}, that will be used frequently and without
explicit reference.) Although this is not posible in each step, we (essentially)
show that for any $E$ there is a polynomial-length sequence of mix operations after which such a
far-apart pair will exist, yielding an overall polynomial bound on a perfect-mixing sequence.
Additionally, to obtain such a perfect-mixing sequence,
at each step we must mix pairs that are either $(\lambda)$-safe, 
for corresponding $\lambda\in\braced{I,I'}$, or near-final.

In this section, for better clarity, we chose to use a top-down presentation style, starting with the statement
of the main theorem (Theorem~\ref{thm: perfect mixability polynomial} below) and its proof, while
all the necessary lemmas that cover different cases from that proof are presented later.
Also, throughout this section, $\gamma$ denotes a small integral constant; for
concreteness we can assume that $\gamma = 2$.

%%%%%%%%%%%%%%%%%

\begin{theorem} \label{thm: perfect mixability polynomial}
	$E$ can be perfectly mixed with precision $0$ 
	in a polynomial number of mixing operations.
\end{theorem} 

\begin{proof}
	If $n = 4$ (and, more generally, if $n$ is a power of $2$), we can simply use 
	Theorem~\ref{them: n power of two} and we are done.
	Thus, we can assume that $n\geq 5$.	
	If $n<22$, we can perfectly mix $E$ using Theorem~\ref{thm: perfect mixability upper bound},
	where the bound on the length of the mixing sequence is polynomial if $n$ is constant.
	Thus in the rest of the argument below we assume that $n \geq 22$.
	
	We show that as long as $E$ satisfies Invariant~(I) (which is the invariant that applies when $n\ge 22$),
	we can find a polynomial-length sequence of mixing operations that will
	decrease $\potential(E)$ by a factor of $1-\Omega(1/n)$. This will be sufficient to prove the theorem.
	
	So assume that $E$ is not yet perfectly mixed and let $\delta=\diameter(E)$. Let $\pi$ be the parity
	for which $\nofdroplets{E_{\pi}}\geq\nofdroplets{E_{\barpi}}$.
	(Recall that we say that $A\subseteq E$ is (I)-mixed if there is no (I)-safe pair in $A$,
	that is, there no pair of distinct concentrations in $A$ whose mixing preserves Invariant~(I) for $E$.
	We will often use this terminology when mixing pairs of concentrations in $E_\pi,E_\barpi\subseteq E$.)
	We consider several cases.
	
	%%%%%%%%%%
	
	\medskip\noindent
	\mycase{1} $\nofconcentrations{E_{\pi}}=1$.
	Let $a\in E_{\pi}$ (that is, $a$ is the only concentration in $E_{\pi}$). 
	Observe that $E_{\barpi}\neq\emptyset$. We have two sub-cases:
	
	\begin{description}

		\item\mycase{1.1} $\min(E_{\bar{\pi}}) < a < \max(E_{\bar{\pi}})$.
		In this case, per Lemma~\ref{lem: farm mix in E odd overlapping E even},
		there are at most two consecutive mixing operations 
		(involving either (I)-safe or near-final pairs)
		after which $\potential(E)$ decreases at least by a factor of $1 - 1/32n$.

		\item \mycase{1.2} $a < \min(E_{\bar{\pi}})$ or $a > \max(E_{\bar{\pi}})$. In this case,
		per Lemma~\ref{lem: far mix in E odd non-overlapping E even}, there is a sequence of
		at most $128ns(E)+1$ (I)-safe mixing operations after which
		$\potential(E)$ decreases at least by a factor of $1 - 1/32\gamma^2n$.
		
	\end{description}	

	%%%%%%%%%%
	
	\medskip\noindent
	\mycase{2} $\nofconcentrations{E_{\pi}}\geq 2$ and $\diameter(E_\pi)\geq\delta/\gamma$.
	In this case, per Lemma~\ref{lem: far mix in E even},	
	there is an (I)-safe pair in $E_{\pi}$ whose mixing decreases
	$\potential(E)$ at least by a factor of $1 - 1/8\gamma^2n$.
	
	%%%%%%%%%%
	
	\medskip\noindent
	\mycase{3} $\nofconcentrations{E_{\pi}}\geq 2$ and $\diameter(E_\pi)<\delta/\gamma$.
	Using Lemma~\ref{lem: E-mix c_pi n >= 22}, there is a sequence of at most
	$128ns(E)$ (I)-safe mixing operations in $E_\pi$, such that after this sequence
	$E_{\pi}$ is (I)-mixed. Let $E'$ be the configuration obtained from $E$
	after applying this sequence. We consider two sub-cases:	
	
	\begin{description}

		\item\mycase{3.1} $\nofdroplets{E'_{\pi}}\geq\nofdroplets{E'_{\bar{\pi}}}$.
		By Observation~\ref{obs: number of same-parity unsafe pairs} below and
		since $E'_{\pi}$ is (I)-mixed, we have that $\nofconcentrations{E'_{\pi}}=1$.		
		Hence, as in Case~$1$ above, there is a sequence of at most $128ns(E')+1$ mixing operations 
		(involving either (I)-safe or near-final pairs)
		after which $\potential(E')$ decreases at least by a factor of $1-1/32\gamma^2n$.
		Thus, $\potential(E)$ also decreases at least by a factor of $1 - 1/32\gamma^2n$.

		\item\mycase{3.2} $\nofdroplets{E'_{\pi}}<\nofdroplets{E'_{\bar{\pi}}}$.
		As $\gamma=2$, $\diameter(E'_\pi)<\delta/\gamma$, and 
		because	the mixing operations on $E_{\pi}$ produced at least one droplet 
		with parity $\bar{\pi}$, we have that 
		$\diameter(E'_{\bar{\pi}})>\delta/\gamma$.
		Hence, as in Case~$2$ above, 
		there is an (I)-safe pair whose mixing decreases
		$\potential(E')$ at least by a factor of $1 - 1/8\gamma^2n$.
		Thus, such mixing also decreases $\potential(E)$ at least by a factor of $1 - 1/8\gamma^2n$.				
	\end{description}
	
	Applying a sequence of mixing operations specified by the cases above
	results in a decrease of $\potential(E)$ at least by a factor of $1 - 1/32\gamma^2n$.
	Thus, by a simple extension to Observation~\ref{obs: far-appart constant mix},
	we obtain that after at most $32\gamma^2n$ such mixing sequences $\potential(E)$ decreases at least by half.
	It follows that after at most $32\gamma^2n\log\potential(E)$ of these mixing sequences (where $\potential(E)$
	denotes the initial potential value), $E$ becomes (I)-mixed, and,
	as $E$ satisfies Invariant~(I), there is a near-final pair in $E$ that we then mix to make $E$ near-final.
	Since $\log\potential(E)\le 2s(E)$, the length of this sequence is at most $64\gamma^2ns(E)$ (where $s(E)$ 
	represents the original size of $E$).
	Each such mixing sequence involves at most $256ns(E)+1$ mixing operations and,
	by~Theorem~\ref{thm: E near-final}, if $E$ is near-final then it can be perfectly mixed 
	by a sequence of at most $144n^3s^2(E)$	mixing operations. 
	Therefore the total number of mixing operations to perfectly mix $E$ is
	at most $2^{14}\gamma^2n^2s^2(E) + 64\gamma^2ns(E) + 144n^3s^2(E)$.
\end{proof}

%%%%%%%%%%%%%%%%%%%%

\begin{observation}\label{obs: number of same-parity unsafe pairs}
	Assume that $E$ with $\nofdroplets{E}=n\geq 22$ satisfies Invariant~(I)
	and let $\pi$ be the parity for which $\nofdroplets{E_{\pi}}\geq\nofdroplets{E_{\bar{\pi}}}$.
	There is at least one droplet in $E_{\pi}$ such that, when paired with any other
	droplet in $E_{\pi}$, the resulting pair is (I)-safe.
\end{observation}

\begin{proof}
	Lemma~$\ref{lem: at most one mod-unsafe pairs}$(a) and the number of distinct odd 
	prime factors of $n$ being less than $\log_3{n}$ imply that $E$ has at most
	$2\lfloor\log_3{n}\rfloor$ droplets that are $\barp$-unsafe when paired with other droplets in $E$.
	This, and Observation~$\ref{obs: decrease number of non-singletons}$ below
	give that the number of droplets that, when mixed with other droplets in $E$
	violate Invariant~(I), is at most 
	$2\lfloor\log_3{n}\rfloor+6< n/2\leq\nofdroplets{E_{\pi}}$, since $n\geq22$.
	Thus, the lemma holds.
\end{proof}
%%%%%%%%%%%%%%%%%%%%

\begin{observation}\label{obs: decrease number of non-singletons}
Assume that $E$ with $\nofdroplets{E}=n\geq 7$ satisfies Invariant~(I).
The number of droplets involved in mixing operations that decrease the number of
non-singletons in $E$ down to one is at most $6$.
\end{observation}

\begin{proof}
	First of all, if the number of non-singletons in $E$ is more than three, 
	then no mixing decreases the number of non-singletons down to one;
	similarly when mixing a non-singleton with frequency higher than two.
	Additionally, mixing two singletons does not decrease the number of non-singletons.	
	Now, $E$ satisfying Invariant~(I) implies that there are 
	at least two non-singletons $a,b\in E$.	
	We consider two types of situations where a mixing decreases the number 
	of non-singletons in $E$ down to one:
	
	\medskip\noindent\mycase{1} Mixing two non-singletons, say $a$ and $b$.
	This can happen when the frequency of both $a$ and $b$ is two each,
	leading to a total of $4$ droplets involved.
	(There could be another non-singleton $e=\half(a+b)$ in $E$, however no mixing 
	involving $e$ would decrease the number of non-singletons down to one
	because either $a$ or $b$ would remain non-singleton after the mixing.)	
	
	\medskip\noindent\mycase{2} Mixing a non-singleton with a singleton.
	(This can happen when $E$ has exactly two non-singletons, $a$ and $b$ respectively.)
	Let $a$ and $c$ be the non-singleton and singleton, respectively.	
	If $a$ is a doubleton and $b=\half(a+c)$, then mixing $a$ and $c$ decreases
	the number of non-singletons down to one.
	Similarly, if $b$ is a doubleton and $a=\half(b+d)$, for some singleton $d\in E$,
	then mixing $b$ and $d$ decreases the number of non-singletons down to one.
	Therefore, the total number of droplets (excluding $a$ and $b$, which were already
	counted in Case~$1$ above) is at most $2$.
	
	\medskip
	Therefore, the number of droplets involved in mixing operations that decrease 
	the number of non-singletons down to one is at most $6$.
\end{proof}

%%%%%%%%%%%%%%%%%%%%

\begin{lemma}\label{lem: farm mix in E odd overlapping E even}
	Assume that $E$ with $\nofdroplets{E}=n\geq 7$ satisfies Invariant~(I).
	Also, assume that $\nofdroplets{E_{\pi}}\geq\nofdroplets{E_{\bar{\pi}}}\geq 1$ and
	that $\nofconcentrations{E_{\pi}} = 1$, with $a\in E_{\pi}$.
	If $\min(E_{\bar{\pi}})<a<\max(E_{\bar{\pi}})$ then there exist at most
	two consecutive mixing operations (involving either (I)-safe or near-final pairs) that
	decrease $\potential(E)$ at least by a factor of $1-1/32n$.
\end{lemma}   

\begin{proof}
	Let $\delta=\diameter(E)$ before any mixing operation.
	Assume without loss of generality that $\pi=even$.
	Since $E$ satisfies Invariant~(I), there is at least one non-singleton $b\in E_{odd}$.
	Either $\min(E_{odd})$ or $\max(E_{odd})$ is furthest from $b$, so
	assume without loss of generality that $\min(E_{odd})$ is furthest from $b$; 
	$|b-\min(E_{odd})|\geq\delta/2$.	
	If pair $(b,\min(E_{odd}))$ is (I)-safe, then we are done.
	So, assume otherwise, namely that pair $(b,\min(E_{odd}))$ is not (I)-safe;
	$a=\half(b+\min(E_{odd}))$ and $b$ (which is a doubleton) are non-singletons in $E$.
	(There could be at most one more non-singleton in $E$, namely $\min(E_{odd})$ strictly doubleton;
	otherwise $(b,\min(E_{odd}))$ would be in fact (I)-safe.)
	Consider the following cases:
	
	\medskip\noindent\mycase{1} $b = \max(E_{odd})$.
	We further consider the following two sub-cases based on the number
	of distinct concentrations in $E$:
	\begin{description}
		\item\mycase{1.1} $\nofconcentrations{E}=3$.
		Then $E=\braced{\nofdroplets{E_{odd}}-2:\min(E_{odd}),\nofdroplets{E_{even}}:a,2:b}$,
		with $\nofdroplets{E_{odd}}\leq 4$.
		If $\nofdroplets{E_{odd}}=3$, then $\min(E_{odd})$ is a singleton, and, 
		after the mixing, $\nofconcentrations{E}=2$ holds,
		which implies that $E$ is now near-final, by Lemma~\ref{lem: mixability m = 2}
		and thus $(b,\min(E_{odd}))$ is actually a near-final pair.		
		Otherwise, $\nofdroplets{E_{odd}}=4$ and $\min(E_{odd})$ is a doubleton.
		This trivially implies that after (and before) the mixing $E$ is near-final,
		so $(b,\min(E_{odd}))$ is a near-final pair.
		(Recall that $|b-\min(E_{odd})|\geq\delta/2$, so the lemma holds
		by Observation~\ref{obs: far-appart mix}.)
				
		
		\item\mycase{1.2} $\nofconcentrations{E}\geq 4$.
		This implies that there is $c\in E_{odd}$ such that $\min(E_{odd})<c<b$;
		mixing $b$ and $c$ produces a non-singleton other than $a$,
		which preserves Invariant~(I), so $(b,c)$ is (I)-safe.
		We analyze the following sub-cases:
		\begin{description}
			\vspace{-0.03in}
			\item\mycase{1.2.1} $c<a$.
			Since $|b-c|\geq\delta/2$, mixing $b$ and $c$ decreases $\potential(E)$ 
			at least by a factor of $1-1/8n$, so we mix them.
			
			\item\mycase{1.2.2} $c>a$.
			Let $d=\half(b+c)$ be the output of the mixing between $b$ and $c$.
			If $d$ is odd, then mixing $d$ and $\min(E_{odd})$ decreases $\potential(E)$ 
			at least by a factor of $1-1/8n$. (Pair $(d, \min(E_{odd}))$ is (I)-safe
			because its mixing produces a non-singleton other than $a$.)
			Otherwise, $d$ is even and, as $|d-a|\geq\delta/4$, 
			mixing $a$ and $d$ decreases $\potential(E)$ at least by a factor of $1-1/32n$.
			(Note that $(a,d)$ is (I)-safe because, after its mixing, $a$ remains non-singleton.)
		\end{description}
		
	\end{description}

	\noindent\mycase{2} $b < \max(E_{odd})$.
	This is similar to Case~$1.2.2$ above under the assumption that $b < \max(E_{odd})$
	and using $c=\max(E_{odd})$.
\end{proof}

%%%%%%%%%%%%%%%%%%%%

\begin{lemma}\label{lem: far mix in E odd non-overlapping E even}
	Assume that $E$ with $\nofdroplets{E}=n\geq 7$ satisfies Invariant~(I).
	Also, assume that $\nofdroplets{E_{\pi}}\geq\nofdroplets{E_{\bar{\pi}}}\geq1$ and
	that $\nofconcentrations{E_\pi} = 1$ with $a\in E_\pi$.
	If either $a < \min(E_{\bar{\pi}})$ or $a>\max(E_{\bar{\pi}})$,	then there exists
	a mixing sequence (of (I)-safe pairs) of length at most $128ns(E) + 1$
	that decreases $\potential(E)$ at least by a factor of $1-1/32\gamma^2n$, 
	for constant $\gamma\in\posintegers$.
\end{lemma}

\begin{proof}
	Let $\delta=\diameter(E)$ before any mixing operation.
	Assume without loss of generality that $\pi=even$ and that $a<\min(E_{odd})$.
	We analyze two sub-cases:
	
	\medskip\noindent\mycase{1} $\diameter(E_{odd})\geq\delta/2\gamma$.
	$E$ satisfying Invariant~(I) implies that there is a non-singleton $b\in E_{odd}$.
	Either $\min(E_{odd})$ or $\max(E_{odd})$ is furthest from $b$, so
	assume without loss of generality that $\min(E_{odd})$ is furthest from $b$.
	As $a<\half(b+\min(E_{odd}))$, pair $(b,\min(E_{odd}))$ is (I)-safe 
	(mixing $b$ and $\min(E_{odd})$	produces a non-singleton other than $a$)
	that satisfies $|b-\min(E_{odd})|\geq\delta/4\gamma$.
	Therefore, mixing $b$ and $\min(E_{odd})$ decreases $\potential(E)$ 
	at least by a factor of $1-1/32\gamma^2n$.
	
	\medskip\noindent\mycase{2} $\diameter(E_{odd})<\delta/2\gamma$.
	This implies that $|\min(E_{odd})-a|\geq\delta/2\gamma$.
	(I)-mix $E_{odd}$ using Lemma~$\ref{lem: E-mix c_pi n >= 22}$;
	$E$ satisfying Invariant~(I) (and our proof in 
	Section~\ref{subsec: proof for n greater-equal than 7}) 
	implies that an even concentration $b$ is eventually produced.
	Since $a$'s frequency is at least $\ceil{n/2}\geq 3$, pair $(a,b)$ is (I)-safe;
	$a$ remains non-singleton after the mixing.
	Additionally, since $b > \min(E_{odd})$, $|b-a|>\delta/2\gamma$ and thus
	mixing $a$ and $b$ decreases $\potential(E)$ at least by a factor of $1-1/8\gamma^2n$.
	Finally, Lemma~$\ref{lem: E-mix c_pi n >= 22}$ takes at most 
	$128ns(E)$ mixing operations, so the lemma holds.
\end{proof}

%%%%%%%%%%%%%%%%%%%%

\begin{lemma}\label{lem: far mix in E even}
	Assume that $E$ with $\nofdroplets{E}=n\geq 22$ satisfies Invariant~$(I)$.
	Also, assume that $\nofdroplets{E_{\pi}}\geq\nofdroplets{E_{\bar{\pi}}}$
	and that $\nofconcentrations{E_{\pi}}\geq 2$.
	If $\diameter(E_\pi) \geq \diameter(E)/\gamma$, for constant $\gamma\in\posintegers$, 
	then there exists an (I)-safe pair whose mixing decreases $\potential(E)$ 
	at least by a factor of $1-1/8\gamma^2n$.
\end{lemma}

\begin{proof}
	Assume without loss of generality that $\pi=even$ and let $a=\min(E_{even})$ and $b=\max(E_{even})$.
	Mixing $a$ and $b$ decreases $\potential(E)$ at least by a factor of $1-1/2\gamma^2n$.
	So, if $(a,b)$ is (I)-safe then we are done.
	Instead, assume otherwise; namely, that $(a,b)$ is not (I)-safe.
	
	The above assumption, and Observation~$\ref{obs: number of same-parity unsafe pairs}$,
	imply that there is $c\in E_{even}$ with $c\notin\braced{a,b}$ 
	for which any pair involving $c$ is (I)-safe.
	Now, either $a$ or $b$ is furthest from $c$, so	assume without loss of generality 
	that $a$ is furthest from $c$.
	Then, $|c-a|\geq\diameter(E)/2\gamma$ and thus mixing $a$ and $c$ decreases 
	$\potential(E)$ at least by a factor of $1-1/8\gamma^2n$.
\end{proof}

%%%%%%%%%%%%%%%%%%%%

\begin{lemma}\label{lem: E-mix c_pi n >= 22}
	Assume that $E$ with $\nofdroplets{E}=n\geq 7$ satisfies Invariant~$(I)$.
	Then, there exists a mixing sequence (of (I)-safe pairs) 
	of length at most $128ns(E)$ that (I)-mixes $E_\pi$.
\end{lemma}

\begin{proof}
	Assume without loss of generality that $\pi = even$ and let $\nofdroplets{E_{even}}=n'$ and
	$\delta=\diameter(E_{even})$ before any mixing operation.
	If a mixing in $E_{even}$ produces droplets with an odd concentration,
	then these droplets are excluded from $E_{even}$ and included in $E_{odd}$;
	this can happen at most $n'/2$ times before $E_{even}$ becomes (I)-mixed.
	
	Assume that the number of (I)-safe pairs in $E_{even}$ is non-zero and 
	let $a=\min(E_{even})$ and $b=\max(E_{even})$.
	If an (I)-safe pair $x,y\in E_{even}$ satisfies $|x-y|\geq \delta/4$,
	then mixing $x$ and $y$ decreases 
	$\potential(E_{even})$ at least by a factor of $1-1/32n'$.
	It follows from Observation~$\ref{obs: far-appart constant mix}$
	that after at most $32n'$ such mixing operations, $\potential(E_{even})$
	decreases at least by a factor of $\half$.
	Hence, $E_{even}$ can be (I)-mixed after at most $32n'\log{\potential(E_{even})}$
	such mixing operations.
	
	We next show that if $E_{even}$ has not been (I)-mixed, then
	after at most two consecutive (I)-safe mixing operations
	either an odd concentration is produced 
	or $\potential(E_{even})$ decreases at least by a factor of $1-1/32n'$.
	Consequently, after at most $128ns(E)$ (I)-safe mixing operations, $E_{even}$ becomes (I)-mixed.
	
	So, if $(a,b)$ is (I)-safe then we are done; $|a-b|= \delta$
	and thus mixing $a$ and $b$ decreases $\potential(E)$ at least by a factor of $1-1/2n$.
	Instead, assume that $(a,b)$ is not (I)-safe and consider the following cases:
	
	\medskip\noindent\mycase{1} $E_{even}$ has only singletons.
	This implies that no mixing in $E_{even}$ decreases the number of non-singletons in $E$.
	So, let $c,d\in E_{even}$ be the furthest-apart (I)-safe pair
	and $x=\half(c+d)$ be the output of their mixing.
	If either $|c-d|\geq \delta/2$ or $x$ odd holds, then we are done.
	Otherwise, let $y\in\braced{a,b}$ be furthest from $x$.
	By the choice of $y$, we have that $|x-y|\geq \delta/2$,
	and as $x$ is a doubleton, $(x,y)$ is (I)-safe and thus we mix it.
	
	\medskip\noindent\mycase{2} $E_{even}$ has exactly one non-singleton $c$.
	Either $a$ or $b$ is furthest from $c$, so assume without loss of generality 
	that $b$ is furthest from $c$; $|b-c|\geq \delta/2$.
	If $(b,c)$ is (I)-safe then we mix it and we are done.
	Otherwise, $(b,c)$ is not (I)-safe, therefore there are exactly two non-singletons in $E$,
	$c\in E_{even}$ (which is a doubleton) and $e=(b+c)/2\notin E_{even}$.
	Since there is an (I)-safe pair in $E_{even}$, there is some $d\in E_{even}$
	such that $d\notin\braced{b,c}$. Then, $(c,d)$ is (I)-safe because 
	mixing $c$ and $d$ produces a non-singleton $x=\half(c+d)$ other than $e$.
	So, mix $c$ and $d$. If $|c-d|\geq\delta/2$ then we are done; similarly if $x$ is odd.
	Otherwise, $(b,x)$ is (I)-safe ($e\neq\half(b+x)$) and $|x-b|\geq\delta/4$, so we mix it.
	
	\medskip\noindent\mycase{3} $E_{even}$ has at least two non-singletons $c<d$.
	We analyze two sub-cases:
	\begin{description}
		\item{\mycase{3.1}} $a=c$.
		Since $(a,b)$ is not (I)-safe, then mixing $b$ and $c$ decreases
		the number of non-singletons down to one. This gives us two sub-cases:
		\begin{description}
			\item{\mycase{3.1.1}} $d=b$.
			$E_{even}$ having an (I)-safe pair implies that there is some $x\in E_{even}$ such that
			$x\notin\braced{c,d}$. Either $c$ or $d$ is furthest from $x$, so			
			assume without loss of generality that $d$ is furthest from $x$; $|d-x|\geq\delta/2$.
			Mixing $x$ and $d$ produces a non-singleton other than $c$, so $(x,d)$ is (I)-safe
			and we mix it.			

			\item{\mycase{3.1.2}} $d<b$.
			This implies that $d=\half(c+b)$, $b-d=\delta/2$ and $(b,d)$ is (I)-safe
			(mixing $b$ and $d$ produces a non-singleton other than $c$), so we mix $b$ and $d$.
		\end{description}

		\item{\mycase{3.2}} $a<c$.
		If $|a-c|\geq\delta/2$ then we mix $a$ and $c$ and we are done;
		$\half(a+c) < d$ implies that mixing $a$ and $c$ produces a non-singleton other than $d$,
		so $(a,c)$ is (I)-safe.
		Instead, let $|a-c|<\delta/2$ which implies $|b-c|\geq\delta/2$.
		Now, if pair $(b,c)$ is (I)-safe then we mix it and we are done, 
		so also assume that $(b,c)$ is not (I)-safe.
		(We cannot have $b=d$ because that would contradict $(a,b)$ not being (I)-safe;
		$|a-c|<\delta/2$ implies $c < \half(a+b)$, so mixing $a$ and $b$ would produce
		a non-singleton other than $c$.)
		Therefore, as $(b,c)$ is not (I)-safe, we have that $d = \half(b+c)$ and mixing $b$ and $d$ 
		produces a non-singleton other than $c$. Hence, $(b,d)$ is (I)-safe
		satisfying $|b-d|\geq\delta/4$, and we mix it.
	\end{description}

	All the pairs mentioned above have difference at least $\delta/4$, so their mixing decreases 
	$\potential(E_{even})$	at least by a factor of $1-1/32n'$.

\end{proof}





We now complete the proof of Theorem~\ref{thm: miscibility characterization}(b).
In Section~\ref{sec: sufficiency of condition mc} we have already established an existence of a perfect-mixing 
sequence for $C$ with intermediate precision $1$.
It still remains to show that this mixing sequence can be modified to have length that is polynomial in $s(C)$.

As explained in the beginning of Section~\ref{sec: polynomial proof}, it remains to
give such a bound for the configuration $E$ constructed from $C$ in Section~\ref{subsec: proof for n greater-equal than 7}
and for mixing sequences with precision $0$ (that is, with only integral concentrations).
Recall that in Section~\ref{sec: polynomial - near final} we gave a polynomial bound if $E$ is near-final.
In Section~\ref{sec: an exponential bound} we gave a bound for arbitrary configurations $E$ that is
exponential only in $n$, thus yielding a polynomial bound for constant $n$.
We will be using these properties in our construction in this section. 
If $E$ is not near-final and $n$ is arbitrary, the general idea for achieving a polynomial bound
is similar to that from Sections~\ref{sec: polynomial - near final} and~\ref{sec: an exponential bound}:
we try to mix a pair of concentrations that are far apart,
to guarantee that we make a quick progress, as measured by the decrease of $\potential(E)$. 
(This is captured by Observation~\ref{obs: far-appart mix}, that will be used frequently and without
explicit reference.) Although this is not posible in each step, we (essentially)
show that for any $E$ there is a polynomial-length sequence of mix operations after which such a
far-apart pair will exist, yielding an overall polynomial bound on a perfect-mixing sequence.
Additionally, to obtain such a perfect-mixing sequence,
at each step we must mix pairs that are either $(\lambda)$-safe, 
for corresponding $\lambda\in\braced{I,I'}$, or near-final.

In this section, for better clarity, we chose to use a top-down presentation style, starting with the statement
of the main theorem (Theorem~\ref{thm: perfect mixability polynomial} below) and its proof, while
all the necessary lemmas that cover different cases from that proof are presented later.
Also, throughout this section, $\gamma$ denotes a small integral constant; for
concreteness we can assume that $\gamma = 2$.

%%%%%%%%%%%%%%%%%

\begin{theorem} \label{thm: perfect mixability polynomial}
	$E$ can be perfectly mixed with precision $0$ 
	in a polynomial number of mixing operations.
\end{theorem} 

\begin{proof}
	If $n = 4$ (and, more generally, if $n$ is a power of $2$), we can simply use 
	Theorem~\ref{them: n power of two} and we are done.
	Thus, we can assume that $n\geq 5$.	
	If $n<22$, we can perfectly mix $E$ using Theorem~\ref{thm: perfect mixability upper bound},
	where the bound on the length of the mixing sequence is polynomial if $n$ is constant.
	Thus in the rest of the argument below we assume that $n \geq 22$.
	
	We show that as long as $E$ satisfies Invariant~(I) (which is the invariant that applies when $n\ge 22$),
	we can find a polynomial-length sequence of mixing operations that will
	decrease $\potential(E)$ by a factor of $1-\Omega(1/n)$. This will be sufficient to prove the theorem.
	
	So assume that $E$ is not yet perfectly mixed and let $\delta=\diameter(E)$. Let $\pi$ be the parity
	for which $\nofdroplets{E_{\pi}}\geq\nofdroplets{E_{\barpi}}$.
	(Recall that we say that $A\subseteq E$ is (I)-mixed if there is no (I)-safe pair in $A$,
	that is, there no pair of distinct concentrations in $A$ whose mixing preserves Invariant~(I) for $E$.
	We will often use this terminology when mixing pairs of concentrations in $E_\pi,E_\barpi\subseteq E$.)
	We consider several cases.
	
	%%%%%%%%%%
	
	\medskip\noindent
	\mycase{1} $\nofconcentrations{E_{\pi}}=1$.
	Let $a\in E_{\pi}$ (that is, $a$ is the only concentration in $E_{\pi}$). 
	Observe that $E_{\barpi}\neq\emptyset$. We have two sub-cases:
	
	\begin{description}

		\item\mycase{1.1} $\min(E_{\bar{\pi}}) < a < \max(E_{\bar{\pi}})$.
		In this case, per Lemma~\ref{lem: farm mix in E odd overlapping E even},
		there are at most two consecutive mixing operations 
		(involving either (I)-safe or near-final pairs)
		after which $\potential(E)$ decreases at least by a factor of $1 - 1/32n$.

		\item \mycase{1.2} $a < \min(E_{\bar{\pi}})$ or $a > \max(E_{\bar{\pi}})$. In this case,
		per Lemma~\ref{lem: far mix in E odd non-overlapping E even}, there is a sequence of
		at most $128ns(E)+1$ (I)-safe mixing operations after which
		$\potential(E)$ decreases at least by a factor of $1 - 1/32\gamma^2n$.
		
	\end{description}	

	%%%%%%%%%%
	
	\medskip\noindent
	\mycase{2} $\nofconcentrations{E_{\pi}}\geq 2$ and $\diameter(E_\pi)\geq\delta/\gamma$.
	In this case, per Lemma~\ref{lem: far mix in E even},	
	there is an (I)-safe pair in $E_{\pi}$ whose mixing decreases
	$\potential(E)$ at least by a factor of $1 - 1/8\gamma^2n$.
	
	%%%%%%%%%%
	
	\medskip\noindent
	\mycase{3} $\nofconcentrations{E_{\pi}}\geq 2$ and $\diameter(E_\pi)<\delta/\gamma$.
	Using Lemma~\ref{lem: E-mix c_pi n >= 22}, there is a sequence of at most
	$128ns(E)$ (I)-safe mixing operations in $E_\pi$, such that after this sequence
	$E_{\pi}$ is (I)-mixed. Let $E'$ be the configuration obtained from $E$
	after applying this sequence. We consider two sub-cases:	
	
	\begin{description}

		\item\mycase{3.1} $\nofdroplets{E'_{\pi}}\geq\nofdroplets{E'_{\bar{\pi}}}$.
		By Observation~\ref{obs: number of same-parity unsafe pairs} below and
		since $E'_{\pi}$ is (I)-mixed, we have that $\nofconcentrations{E'_{\pi}}=1$.		
		Hence, as in Case~$1$ above, there is a sequence of at most $128ns(E')+1$ mixing operations 
		(involving either (I)-safe or near-final pairs)
		after which $\potential(E')$ decreases at least by a factor of $1-1/32\gamma^2n$.
		Thus, $\potential(E)$ also decreases at least by a factor of $1 - 1/32\gamma^2n$.

		\item\mycase{3.2} $\nofdroplets{E'_{\pi}}<\nofdroplets{E'_{\bar{\pi}}}$.
		As $\gamma=2$, $\diameter(E'_\pi)<\delta/\gamma$, and 
		because	the mixing operations on $E_{\pi}$ produced at least one droplet 
		with parity $\bar{\pi}$, we have that 
		$\diameter(E'_{\bar{\pi}})>\delta/\gamma$.
		Hence, as in Case~$2$ above, 
		there is an (I)-safe pair whose mixing decreases
		$\potential(E')$ at least by a factor of $1 - 1/8\gamma^2n$.
		Thus, such mixing also decreases $\potential(E)$ at least by a factor of $1 - 1/8\gamma^2n$.				
	\end{description}
	
	Applying a sequence of mixing operations specified by the cases above
	results in a decrease of $\potential(E)$ at least by a factor of $1 - 1/32\gamma^2n$.
	Thus, by a simple extension to Observation~\ref{obs: far-appart constant mix},
	we obtain that after at most $32\gamma^2n$ such mixing sequences $\potential(E)$ decreases at least by half.
	It follows that after at most $32\gamma^2n\log\potential(E)$ of these mixing sequences (where $\potential(E)$
	denotes the initial potential value), $E$ becomes (I)-mixed, and,
	as $E$ satisfies Invariant~(I), there is a near-final pair in $E$ that we then mix to make $E$ near-final.
	Since $\log\potential(E)\le 2s(E)$, the length of this sequence is at most $64\gamma^2ns(E)$ (where $s(E)$ 
	represents the original size of $E$).
	Each such mixing sequence involves at most $256ns(E)+1$ mixing operations and,
	by~Theorem~\ref{thm: E near-final}, if $E$ is near-final then it can be perfectly mixed 
	by a sequence of at most $144n^3s^2(E)$	mixing operations. 
	Therefore the total number of mixing operations to perfectly mix $E$ is
	at most $2^{14}\gamma^2n^2s^2(E) + 64\gamma^2ns(E) + 144n^3s^2(E)$.
\end{proof}

%%%%%%%%%%%%%%%%%%%%

\begin{observation}\label{obs: number of same-parity unsafe pairs}
	Assume that $E$ with $\nofdroplets{E}=n\geq 22$ satisfies Invariant~(I)
	and let $\pi$ be the parity for which $\nofdroplets{E_{\pi}}\geq\nofdroplets{E_{\bar{\pi}}}$.
	There is at least one droplet in $E_{\pi}$ such that, when paired with any other
	droplet in $E_{\pi}$, the resulting pair is (I)-safe.
\end{observation}

\begin{proof}
	Lemma~$\ref{lem: at most one mod-unsafe pairs}$(a) and the number of distinct odd 
	prime factors of $n$ being less than $\log_3{n}$ imply that $E$ has at most
	$2\lfloor\log_3{n}\rfloor$ droplets that are $\barp$-unsafe when paired with other droplets in $E$.
	This, and Observation~$\ref{obs: decrease number of non-singletons}$ below
	give that the number of droplets that, when mixed with other droplets in $E$
	violate Invariant~(I), is at most 
	$2\lfloor\log_3{n}\rfloor+6< n/2\leq\nofdroplets{E_{\pi}}$, since $n\geq22$.
	Thus, the lemma holds.
\end{proof}
%%%%%%%%%%%%%%%%%%%%

\begin{observation}\label{obs: decrease number of non-singletons}
Assume that $E$ with $\nofdroplets{E}=n\geq 7$ satisfies Invariant~(I).
The number of droplets involved in mixing operations that decrease the number of
non-singletons in $E$ down to one is at most $6$.
\end{observation}

\begin{proof}
	First of all, if the number of non-singletons in $E$ is more than three, 
	then no mixing decreases the number of non-singletons down to one;
	similarly when mixing a non-singleton with frequency higher than two.
	Additionally, mixing two singletons does not decrease the number of non-singletons.	
	Now, $E$ satisfying Invariant~(I) implies that there are 
	at least two non-singletons $a,b\in E$.	
	We consider two types of situations where a mixing decreases the number 
	of non-singletons in $E$ down to one:
	
	\medskip\noindent\mycase{1} Mixing two non-singletons, say $a$ and $b$.
	This can happen when the frequency of both $a$ and $b$ is two each,
	leading to a total of $4$ droplets involved.
	(There could be another non-singleton $e=\half(a+b)$ in $E$, however no mixing 
	involving $e$ would decrease the number of non-singletons down to one
	because either $a$ or $b$ would remain non-singleton after the mixing.)	
	
	\medskip\noindent\mycase{2} Mixing a non-singleton with a singleton.
	(This can happen when $E$ has exactly two non-singletons, $a$ and $b$ respectively.)
	Let $a$ and $c$ be the non-singleton and singleton, respectively.	
	If $a$ is a doubleton and $b=\half(a+c)$, then mixing $a$ and $c$ decreases
	the number of non-singletons down to one.
	Similarly, if $b$ is a doubleton and $a=\half(b+d)$, for some singleton $d\in E$,
	then mixing $b$ and $d$ decreases the number of non-singletons down to one.
	Therefore, the total number of droplets (excluding $a$ and $b$, which were already
	counted in Case~$1$ above) is at most $2$.
	
	\medskip
	Therefore, the number of droplets involved in mixing operations that decrease 
	the number of non-singletons down to one is at most $6$.
\end{proof}

%%%%%%%%%%%%%%%%%%%%

\begin{lemma}\label{lem: farm mix in E odd overlapping E even}
	Assume that $E$ with $\nofdroplets{E}=n\geq 7$ satisfies Invariant~(I).
	Also, assume that $\nofdroplets{E_{\pi}}\geq\nofdroplets{E_{\bar{\pi}}}\geq 1$ and
	that $\nofconcentrations{E_{\pi}} = 1$, with $a\in E_{\pi}$.
	If $\min(E_{\bar{\pi}})<a<\max(E_{\bar{\pi}})$ then there exist at most
	two consecutive mixing operations (involving either (I)-safe or near-final pairs) that
	decrease $\potential(E)$ at least by a factor of $1-1/32n$.
\end{lemma}   

\begin{proof}
	Let $\delta=\diameter(E)$ before any mixing operation.
	Assume without loss of generality that $\pi=even$.
	Since $E$ satisfies Invariant~(I), there is at least one non-singleton $b\in E_{odd}$.
	Either $\min(E_{odd})$ or $\max(E_{odd})$ is furthest from $b$, so
	assume without loss of generality that $\min(E_{odd})$ is furthest from $b$; 
	$|b-\min(E_{odd})|\geq\delta/2$.	
	If pair $(b,\min(E_{odd}))$ is (I)-safe, then we are done.
	So, assume otherwise, namely that pair $(b,\min(E_{odd}))$ is not (I)-safe;
	$a=\half(b+\min(E_{odd}))$ and $b$ (which is a doubleton) are non-singletons in $E$.
	(There could be at most one more non-singleton in $E$, namely $\min(E_{odd})$ strictly doubleton;
	otherwise $(b,\min(E_{odd}))$ would be in fact (I)-safe.)
	Consider the following cases:
	
	\medskip\noindent\mycase{1} $b = \max(E_{odd})$.
	We further consider the following two sub-cases based on the number
	of distinct concentrations in $E$:
	\begin{description}
		\item\mycase{1.1} $\nofconcentrations{E}=3$.
		Then $E=\braced{\nofdroplets{E_{odd}}-2:\min(E_{odd}),\nofdroplets{E_{even}}:a,2:b}$,
		with $\nofdroplets{E_{odd}}\leq 4$.
		If $\nofdroplets{E_{odd}}=3$, then $\min(E_{odd})$ is a singleton, and, 
		after the mixing, $\nofconcentrations{E}=2$ holds,
		which implies that $E$ is now near-final, by Lemma~\ref{lem: mixability m = 2}
		and thus $(b,\min(E_{odd}))$ is actually a near-final pair.		
		Otherwise, $\nofdroplets{E_{odd}}=4$ and $\min(E_{odd})$ is a doubleton.
		This trivially implies that after (and before) the mixing $E$ is near-final,
		so $(b,\min(E_{odd}))$ is a near-final pair.
		(Recall that $|b-\min(E_{odd})|\geq\delta/2$, so the lemma holds
		by Observation~\ref{obs: far-appart mix}.)
				
		
		\item\mycase{1.2} $\nofconcentrations{E}\geq 4$.
		This implies that there is $c\in E_{odd}$ such that $\min(E_{odd})<c<b$;
		mixing $b$ and $c$ produces a non-singleton other than $a$,
		which preserves Invariant~(I), so $(b,c)$ is (I)-safe.
		We analyze the following sub-cases:
		\begin{description}
			\vspace{-0.03in}
			\item\mycase{1.2.1} $c<a$.
			Since $|b-c|\geq\delta/2$, mixing $b$ and $c$ decreases $\potential(E)$ 
			at least by a factor of $1-1/8n$, so we mix them.
			
			\item\mycase{1.2.2} $c>a$.
			Let $d=\half(b+c)$ be the output of the mixing between $b$ and $c$.
			If $d$ is odd, then mixing $d$ and $\min(E_{odd})$ decreases $\potential(E)$ 
			at least by a factor of $1-1/8n$. (Pair $(d, \min(E_{odd}))$ is (I)-safe
			because its mixing produces a non-singleton other than $a$.)
			Otherwise, $d$ is even and, as $|d-a|\geq\delta/4$, 
			mixing $a$ and $d$ decreases $\potential(E)$ at least by a factor of $1-1/32n$.
			(Note that $(a,d)$ is (I)-safe because, after its mixing, $a$ remains non-singleton.)
		\end{description}
		
	\end{description}

	\noindent\mycase{2} $b < \max(E_{odd})$.
	This is similar to Case~$1.2.2$ above under the assumption that $b < \max(E_{odd})$
	and using $c=\max(E_{odd})$.
\end{proof}

%%%%%%%%%%%%%%%%%%%%

\begin{lemma}\label{lem: far mix in E odd non-overlapping E even}
	Assume that $E$ with $\nofdroplets{E}=n\geq 7$ satisfies Invariant~(I).
	Also, assume that $\nofdroplets{E_{\pi}}\geq\nofdroplets{E_{\bar{\pi}}}\geq1$ and
	that $\nofconcentrations{E_\pi} = 1$ with $a\in E_\pi$.
	If either $a < \min(E_{\bar{\pi}})$ or $a>\max(E_{\bar{\pi}})$,	then there exists
	a mixing sequence (of (I)-safe pairs) of length at most $128ns(E) + 1$
	that decreases $\potential(E)$ at least by a factor of $1-1/32\gamma^2n$, 
	for constant $\gamma\in\posintegers$.
\end{lemma}

\begin{proof}
	Let $\delta=\diameter(E)$ before any mixing operation.
	Assume without loss of generality that $\pi=even$ and that $a<\min(E_{odd})$.
	We analyze two sub-cases:
	
	\medskip\noindent\mycase{1} $\diameter(E_{odd})\geq\delta/2\gamma$.
	$E$ satisfying Invariant~(I) implies that there is a non-singleton $b\in E_{odd}$.
	Either $\min(E_{odd})$ or $\max(E_{odd})$ is furthest from $b$, so
	assume without loss of generality that $\min(E_{odd})$ is furthest from $b$.
	As $a<\half(b+\min(E_{odd}))$, pair $(b,\min(E_{odd}))$ is (I)-safe 
	(mixing $b$ and $\min(E_{odd})$	produces a non-singleton other than $a$)
	that satisfies $|b-\min(E_{odd})|\geq\delta/4\gamma$.
	Therefore, mixing $b$ and $\min(E_{odd})$ decreases $\potential(E)$ 
	at least by a factor of $1-1/32\gamma^2n$.
	
	\medskip\noindent\mycase{2} $\diameter(E_{odd})<\delta/2\gamma$.
	This implies that $|\min(E_{odd})-a|\geq\delta/2\gamma$.
	(I)-mix $E_{odd}$ using Lemma~$\ref{lem: E-mix c_pi n >= 22}$;
	$E$ satisfying Invariant~(I) (and our proof in 
	Section~\ref{subsec: proof for n greater-equal than 7}) 
	implies that an even concentration $b$ is eventually produced.
	Since $a$'s frequency is at least $\ceil{n/2}\geq 3$, pair $(a,b)$ is (I)-safe;
	$a$ remains non-singleton after the mixing.
	Additionally, since $b > \min(E_{odd})$, $|b-a|>\delta/2\gamma$ and thus
	mixing $a$ and $b$ decreases $\potential(E)$ at least by a factor of $1-1/8\gamma^2n$.
	Finally, Lemma~$\ref{lem: E-mix c_pi n >= 22}$ takes at most 
	$128ns(E)$ mixing operations, so the lemma holds.
\end{proof}

%%%%%%%%%%%%%%%%%%%%

\begin{lemma}\label{lem: far mix in E even}
	Assume that $E$ with $\nofdroplets{E}=n\geq 22$ satisfies Invariant~$(I)$.
	Also, assume that $\nofdroplets{E_{\pi}}\geq\nofdroplets{E_{\bar{\pi}}}$
	and that $\nofconcentrations{E_{\pi}}\geq 2$.
	If $\diameter(E_\pi) \geq \diameter(E)/\gamma$, for constant $\gamma\in\posintegers$, 
	then there exists an (I)-safe pair whose mixing decreases $\potential(E)$ 
	at least by a factor of $1-1/8\gamma^2n$.
\end{lemma}

\begin{proof}
	Assume without loss of generality that $\pi=even$ and let $a=\min(E_{even})$ and $b=\max(E_{even})$.
	Mixing $a$ and $b$ decreases $\potential(E)$ at least by a factor of $1-1/2\gamma^2n$.
	So, if $(a,b)$ is (I)-safe then we are done.
	Instead, assume otherwise; namely, that $(a,b)$ is not (I)-safe.
	
	The above assumption, and Observation~$\ref{obs: number of same-parity unsafe pairs}$,
	imply that there is $c\in E_{even}$ with $c\notin\braced{a,b}$ 
	for which any pair involving $c$ is (I)-safe.
	Now, either $a$ or $b$ is furthest from $c$, so	assume without loss of generality 
	that $a$ is furthest from $c$.
	Then, $|c-a|\geq\diameter(E)/2\gamma$ and thus mixing $a$ and $c$ decreases 
	$\potential(E)$ at least by a factor of $1-1/8\gamma^2n$.
\end{proof}

%%%%%%%%%%%%%%%%%%%%

\begin{lemma}\label{lem: E-mix c_pi n >= 22}
	Assume that $E$ with $\nofdroplets{E}=n\geq 7$ satisfies Invariant~$(I)$.
	Then, there exists a mixing sequence (of (I)-safe pairs) 
	of length at most $128ns(E)$ that (I)-mixes $E_\pi$.
\end{lemma}

\begin{proof}
	Assume without loss of generality that $\pi = even$ and let $\nofdroplets{E_{even}}=n'$ and
	$\delta=\diameter(E_{even})$ before any mixing operation.
	If a mixing in $E_{even}$ produces droplets with an odd concentration,
	then these droplets are excluded from $E_{even}$ and included in $E_{odd}$;
	this can happen at most $n'/2$ times before $E_{even}$ becomes (I)-mixed.
	
	Assume that the number of (I)-safe pairs in $E_{even}$ is non-zero and 
	let $a=\min(E_{even})$ and $b=\max(E_{even})$.
	If an (I)-safe pair $x,y\in E_{even}$ satisfies $|x-y|\geq \delta/4$,
	then mixing $x$ and $y$ decreases 
	$\potential(E_{even})$ at least by a factor of $1-1/32n'$.
	It follows from Observation~$\ref{obs: far-appart constant mix}$
	that after at most $32n'$ such mixing operations, $\potential(E_{even})$
	decreases at least by a factor of $\half$.
	Hence, $E_{even}$ can be (I)-mixed after at most $32n'\log{\potential(E_{even})}$
	such mixing operations.
	
	We next show that if $E_{even}$ has not been (I)-mixed, then
	after at most two consecutive (I)-safe mixing operations
	either an odd concentration is produced 
	or $\potential(E_{even})$ decreases at least by a factor of $1-1/32n'$.
	Consequently, after at most $128ns(E)$ (I)-safe mixing operations, $E_{even}$ becomes (I)-mixed.
	
	So, if $(a,b)$ is (I)-safe then we are done; $|a-b|= \delta$
	and thus mixing $a$ and $b$ decreases $\potential(E)$ at least by a factor of $1-1/2n$.
	Instead, assume that $(a,b)$ is not (I)-safe and consider the following cases:
	
	\medskip\noindent\mycase{1} $E_{even}$ has only singletons.
	This implies that no mixing in $E_{even}$ decreases the number of non-singletons in $E$.
	So, let $c,d\in E_{even}$ be the furthest-apart (I)-safe pair
	and $x=\half(c+d)$ be the output of their mixing.
	If either $|c-d|\geq \delta/2$ or $x$ odd holds, then we are done.
	Otherwise, let $y\in\braced{a,b}$ be furthest from $x$.
	By the choice of $y$, we have that $|x-y|\geq \delta/2$,
	and as $x$ is a doubleton, $(x,y)$ is (I)-safe and thus we mix it.
	
	\medskip\noindent\mycase{2} $E_{even}$ has exactly one non-singleton $c$.
	Either $a$ or $b$ is furthest from $c$, so assume without loss of generality 
	that $b$ is furthest from $c$; $|b-c|\geq \delta/2$.
	If $(b,c)$ is (I)-safe then we mix it and we are done.
	Otherwise, $(b,c)$ is not (I)-safe, therefore there are exactly two non-singletons in $E$,
	$c\in E_{even}$ (which is a doubleton) and $e=(b+c)/2\notin E_{even}$.
	Since there is an (I)-safe pair in $E_{even}$, there is some $d\in E_{even}$
	such that $d\notin\braced{b,c}$. Then, $(c,d)$ is (I)-safe because 
	mixing $c$ and $d$ produces a non-singleton $x=\half(c+d)$ other than $e$.
	So, mix $c$ and $d$. If $|c-d|\geq\delta/2$ then we are done; similarly if $x$ is odd.
	Otherwise, $(b,x)$ is (I)-safe ($e\neq\half(b+x)$) and $|x-b|\geq\delta/4$, so we mix it.
	
	\medskip\noindent\mycase{3} $E_{even}$ has at least two non-singletons $c<d$.
	We analyze two sub-cases:
	\begin{description}
		\item{\mycase{3.1}} $a=c$.
		Since $(a,b)$ is not (I)-safe, then mixing $b$ and $c$ decreases
		the number of non-singletons down to one. This gives us two sub-cases:
		\begin{description}
			\item{\mycase{3.1.1}} $d=b$.
			$E_{even}$ having an (I)-safe pair implies that there is some $x\in E_{even}$ such that
			$x\notin\braced{c,d}$. Either $c$ or $d$ is furthest from $x$, so			
			assume without loss of generality that $d$ is furthest from $x$; $|d-x|\geq\delta/2$.
			Mixing $x$ and $d$ produces a non-singleton other than $c$, so $(x,d)$ is (I)-safe
			and we mix it.			

			\item{\mycase{3.1.2}} $d<b$.
			This implies that $d=\half(c+b)$, $b-d=\delta/2$ and $(b,d)$ is (I)-safe
			(mixing $b$ and $d$ produces a non-singleton other than $c$), so we mix $b$ and $d$.
		\end{description}

		\item{\mycase{3.2}} $a<c$.
		If $|a-c|\geq\delta/2$ then we mix $a$ and $c$ and we are done;
		$\half(a+c) < d$ implies that mixing $a$ and $c$ produces a non-singleton other than $d$,
		so $(a,c)$ is (I)-safe.
		Instead, let $|a-c|<\delta/2$ which implies $|b-c|\geq\delta/2$.
		Now, if pair $(b,c)$ is (I)-safe then we mix it and we are done, 
		so also assume that $(b,c)$ is not (I)-safe.
		(We cannot have $b=d$ because that would contradict $(a,b)$ not being (I)-safe;
		$|a-c|<\delta/2$ implies $c < \half(a+b)$, so mixing $a$ and $b$ would produce
		a non-singleton other than $c$.)
		Therefore, as $(b,c)$ is not (I)-safe, we have that $d = \half(b+c)$ and mixing $b$ and $d$ 
		produces a non-singleton other than $c$. Hence, $(b,d)$ is (I)-safe
		satisfying $|b-d|\geq\delta/4$, and we mix it.
	\end{description}

	All the pairs mentioned above have difference at least $\delta/4$, so their mixing decreases 
	$\potential(E_{even})$	at least by a factor of $1-1/32n'$.

\end{proof}





We now complete the proof of Theorem~\ref{thm: miscibility characterization}(b).
In Section~\ref{sec: sufficiency of condition mc} we have already established an existence of a perfect-mixing 
sequence for $C$ with intermediate precision $1$.
It still remains to show that this mixing sequence can be modified to have length that is polynomial in $s(C)$.

As explained in the beginning of Section~\ref{sec: polynomial proof}, it remains to
give such a bound for the configuration $E$ constructed from $C$ in Section~\ref{subsec: proof for n greater-equal than 7}
and for mixing sequences with precision $0$ (that is, with only integral concentrations).
Recall that in Section~\ref{sec: polynomial - near final} we gave a polynomial bound if $E$ is near-final.
In Section~\ref{sec: an exponential bound} we gave a bound for arbitrary configurations $E$ that is
exponential only in $n$, thus yielding a polynomial bound for constant $n$.
We will be using these properties in our construction in this section. 
If $E$ is not near-final and $n$ is arbitrary, the general idea for achieving a polynomial bound
is similar to that from Sections~\ref{sec: polynomial - near final} and~\ref{sec: an exponential bound}:
we try to mix a pair of concentrations that are far apart,
to guarantee that we make a quick progress, as measured by the decrease of $\potential(E)$. 
(This is captured by Observation~\ref{obs: far-appart mix}, that will be used frequently and without
explicit reference.) Although this is not posible in each step, we (essentially)
show that for any $E$ there is a polynomial-length sequence of mix operations after which such a
far-apart pair will exist, yielding an overall polynomial bound on a perfect-mixing sequence.
Additionally, to obtain such a perfect-mixing sequence,
at each step we must mix pairs that are either $(\lambda)$-safe, 
for corresponding $\lambda\in\braced{I,I'}$, or near-final.

In this section, for better clarity, we chose to use a top-down presentation style, starting with the statement
of the main theorem (Theorem~\ref{thm: perfect mixability polynomial} below) and its proof, while
all the necessary lemmas that cover different cases from that proof are presented later.
Also, throughout this section, $\gamma$ denotes a small integral constant; for
concreteness we can assume that $\gamma = 2$.

%%%%%%%%%%%%%%%%%

\begin{theorem} \label{thm: perfect mixability polynomial}
	$E$ can be perfectly mixed with precision $0$ 
	in a polynomial number of mixing operations.
\end{theorem} 

\begin{proof}
	If $n = 4$ (and, more generally, if $n$ is a power of $2$), we can simply use 
	Theorem~\ref{them: n power of two} and we are done.
	Thus, we can assume that $n\geq 5$.	
	If $n<22$, we can perfectly mix $E$ using Theorem~\ref{thm: perfect mixability upper bound},
	where the bound on the length of the mixing sequence is polynomial if $n$ is constant.
	Thus in the rest of the argument below we assume that $n \geq 22$.
	
	We show that as long as $E$ satisfies Invariant~(I) (which is the invariant that applies when $n\ge 22$),
	we can find a polynomial-length sequence of mixing operations that will
	decrease $\potential(E)$ by a factor of $1-\Omega(1/n)$. This will be sufficient to prove the theorem.
	
	So assume that $E$ is not yet perfectly mixed and let $\delta=\diameter(E)$. Let $\pi$ be the parity
	for which $\nofdroplets{E_{\pi}}\geq\nofdroplets{E_{\barpi}}$.
	(Recall that we say that $A\subseteq E$ is (I)-mixed if there is no (I)-safe pair in $A$,
	that is, there no pair of distinct concentrations in $A$ whose mixing preserves Invariant~(I) for $E$.
	We will often use this terminology when mixing pairs of concentrations in $E_\pi,E_\barpi\subseteq E$.)
	We consider several cases.
	
	%%%%%%%%%%
	
	\medskip\noindent
	\mycase{1} $\nofconcentrations{E_{\pi}}=1$.
	Let $a\in E_{\pi}$ (that is, $a$ is the only concentration in $E_{\pi}$). 
	Observe that $E_{\barpi}\neq\emptyset$. We have two sub-cases:
	
	\begin{description}

		\item\mycase{1.1} $\min(E_{\bar{\pi}}) < a < \max(E_{\bar{\pi}})$.
		In this case, per Lemma~\ref{lem: farm mix in E odd overlapping E even},
		there are at most two consecutive mixing operations 
		(involving either (I)-safe or near-final pairs)
		after which $\potential(E)$ decreases at least by a factor of $1 - 1/32n$.

		\item \mycase{1.2} $a < \min(E_{\bar{\pi}})$ or $a > \max(E_{\bar{\pi}})$. In this case,
		per Lemma~\ref{lem: far mix in E odd non-overlapping E even}, there is a sequence of
		at most $128ns(E)+1$ (I)-safe mixing operations after which
		$\potential(E)$ decreases at least by a factor of $1 - 1/32\gamma^2n$.
		
	\end{description}	

	%%%%%%%%%%
	
	\medskip\noindent
	\mycase{2} $\nofconcentrations{E_{\pi}}\geq 2$ and $\diameter(E_\pi)\geq\delta/\gamma$.
	In this case, per Lemma~\ref{lem: far mix in E even},	
	there is an (I)-safe pair in $E_{\pi}$ whose mixing decreases
	$\potential(E)$ at least by a factor of $1 - 1/8\gamma^2n$.
	
	%%%%%%%%%%
	
	\medskip\noindent
	\mycase{3} $\nofconcentrations{E_{\pi}}\geq 2$ and $\diameter(E_\pi)<\delta/\gamma$.
	Using Lemma~\ref{lem: E-mix c_pi n >= 22}, there is a sequence of at most
	$128ns(E)$ (I)-safe mixing operations in $E_\pi$, such that after this sequence
	$E_{\pi}$ is (I)-mixed. Let $E'$ be the configuration obtained from $E$
	after applying this sequence. We consider two sub-cases:	
	
	\begin{description}

		\item\mycase{3.1} $\nofdroplets{E'_{\pi}}\geq\nofdroplets{E'_{\bar{\pi}}}$.
		By Observation~\ref{obs: number of same-parity unsafe pairs} below and
		since $E'_{\pi}$ is (I)-mixed, we have that $\nofconcentrations{E'_{\pi}}=1$.		
		Hence, as in Case~$1$ above, there is a sequence of at most $128ns(E')+1$ mixing operations 
		(involving either (I)-safe or near-final pairs)
		after which $\potential(E')$ decreases at least by a factor of $1-1/32\gamma^2n$.
		Thus, $\potential(E)$ also decreases at least by a factor of $1 - 1/32\gamma^2n$.

		\item\mycase{3.2} $\nofdroplets{E'_{\pi}}<\nofdroplets{E'_{\bar{\pi}}}$.
		As $\gamma=2$, $\diameter(E'_\pi)<\delta/\gamma$, and 
		because	the mixing operations on $E_{\pi}$ produced at least one droplet 
		with parity $\bar{\pi}$, we have that 
		$\diameter(E'_{\bar{\pi}})>\delta/\gamma$.
		Hence, as in Case~$2$ above, 
		there is an (I)-safe pair whose mixing decreases
		$\potential(E')$ at least by a factor of $1 - 1/8\gamma^2n$.
		Thus, such mixing also decreases $\potential(E)$ at least by a factor of $1 - 1/8\gamma^2n$.				
	\end{description}
	
	Applying a sequence of mixing operations specified by the cases above
	results in a decrease of $\potential(E)$ at least by a factor of $1 - 1/32\gamma^2n$.
	Thus, by a simple extension to Observation~\ref{obs: far-appart constant mix},
	we obtain that after at most $32\gamma^2n$ such mixing sequences $\potential(E)$ decreases at least by half.
	It follows that after at most $32\gamma^2n\log\potential(E)$ of these mixing sequences (where $\potential(E)$
	denotes the initial potential value), $E$ becomes (I)-mixed, and,
	as $E$ satisfies Invariant~(I), there is a near-final pair in $E$ that we then mix to make $E$ near-final.
	Since $\log\potential(E)\le 2s(E)$, the length of this sequence is at most $64\gamma^2ns(E)$ (where $s(E)$ 
	represents the original size of $E$).
	Each such mixing sequence involves at most $256ns(E)+1$ mixing operations and,
	by~Theorem~\ref{thm: E near-final}, if $E$ is near-final then it can be perfectly mixed 
	by a sequence of at most $144n^3s^2(E)$	mixing operations. 
	Therefore the total number of mixing operations to perfectly mix $E$ is
	at most $2^{14}\gamma^2n^2s^2(E) + 64\gamma^2ns(E) + 144n^3s^2(E)$.
\end{proof}

%%%%%%%%%%%%%%%%%%%%

\begin{observation}\label{obs: number of same-parity unsafe pairs}
	Assume that $E$ with $\nofdroplets{E}=n\geq 22$ satisfies Invariant~(I)
	and let $\pi$ be the parity for which $\nofdroplets{E_{\pi}}\geq\nofdroplets{E_{\bar{\pi}}}$.
	There is at least one droplet in $E_{\pi}$ such that, when paired with any other
	droplet in $E_{\pi}$, the resulting pair is (I)-safe.
\end{observation}

\begin{proof}
	Lemma~$\ref{lem: at most one mod-unsafe pairs}$(a) and the number of distinct odd 
	prime factors of $n$ being less than $\log_3{n}$ imply that $E$ has at most
	$2\lfloor\log_3{n}\rfloor$ droplets that are $\barp$-unsafe when paired with other droplets in $E$.
	This, and Observation~$\ref{obs: decrease number of non-singletons}$ below
	give that the number of droplets that, when mixed with other droplets in $E$
	violate Invariant~(I), is at most 
	$2\lfloor\log_3{n}\rfloor+6< n/2\leq\nofdroplets{E_{\pi}}$, since $n\geq22$.
	Thus, the lemma holds.
\end{proof}
%%%%%%%%%%%%%%%%%%%%

\begin{observation}\label{obs: decrease number of non-singletons}
Assume that $E$ with $\nofdroplets{E}=n\geq 7$ satisfies Invariant~(I).
The number of droplets involved in mixing operations that decrease the number of
non-singletons in $E$ down to one is at most $6$.
\end{observation}

\begin{proof}
	First of all, if the number of non-singletons in $E$ is more than three, 
	then no mixing decreases the number of non-singletons down to one;
	similarly when mixing a non-singleton with frequency higher than two.
	Additionally, mixing two singletons does not decrease the number of non-singletons.	
	Now, $E$ satisfying Invariant~(I) implies that there are 
	at least two non-singletons $a,b\in E$.	
	We consider two types of situations where a mixing decreases the number 
	of non-singletons in $E$ down to one:
	
	\medskip\noindent\mycase{1} Mixing two non-singletons, say $a$ and $b$.
	This can happen when the frequency of both $a$ and $b$ is two each,
	leading to a total of $4$ droplets involved.
	(There could be another non-singleton $e=\half(a+b)$ in $E$, however no mixing 
	involving $e$ would decrease the number of non-singletons down to one
	because either $a$ or $b$ would remain non-singleton after the mixing.)	
	
	\medskip\noindent\mycase{2} Mixing a non-singleton with a singleton.
	(This can happen when $E$ has exactly two non-singletons, $a$ and $b$ respectively.)
	Let $a$ and $c$ be the non-singleton and singleton, respectively.	
	If $a$ is a doubleton and $b=\half(a+c)$, then mixing $a$ and $c$ decreases
	the number of non-singletons down to one.
	Similarly, if $b$ is a doubleton and $a=\half(b+d)$, for some singleton $d\in E$,
	then mixing $b$ and $d$ decreases the number of non-singletons down to one.
	Therefore, the total number of droplets (excluding $a$ and $b$, which were already
	counted in Case~$1$ above) is at most $2$.
	
	\medskip
	Therefore, the number of droplets involved in mixing operations that decrease 
	the number of non-singletons down to one is at most $6$.
\end{proof}

%%%%%%%%%%%%%%%%%%%%

\begin{lemma}\label{lem: farm mix in E odd overlapping E even}
	Assume that $E$ with $\nofdroplets{E}=n\geq 7$ satisfies Invariant~(I).
	Also, assume that $\nofdroplets{E_{\pi}}\geq\nofdroplets{E_{\bar{\pi}}}\geq 1$ and
	that $\nofconcentrations{E_{\pi}} = 1$, with $a\in E_{\pi}$.
	If $\min(E_{\bar{\pi}})<a<\max(E_{\bar{\pi}})$ then there exist at most
	two consecutive mixing operations (involving either (I)-safe or near-final pairs) that
	decrease $\potential(E)$ at least by a factor of $1-1/32n$.
\end{lemma}   

\begin{proof}
	Let $\delta=\diameter(E)$ before any mixing operation.
	Assume without loss of generality that $\pi=even$.
	Since $E$ satisfies Invariant~(I), there is at least one non-singleton $b\in E_{odd}$.
	Either $\min(E_{odd})$ or $\max(E_{odd})$ is furthest from $b$, so
	assume without loss of generality that $\min(E_{odd})$ is furthest from $b$; 
	$|b-\min(E_{odd})|\geq\delta/2$.	
	If pair $(b,\min(E_{odd}))$ is (I)-safe, then we are done.
	So, assume otherwise, namely that pair $(b,\min(E_{odd}))$ is not (I)-safe;
	$a=\half(b+\min(E_{odd}))$ and $b$ (which is a doubleton) are non-singletons in $E$.
	(There could be at most one more non-singleton in $E$, namely $\min(E_{odd})$ strictly doubleton;
	otherwise $(b,\min(E_{odd}))$ would be in fact (I)-safe.)
	Consider the following cases:
	
	\medskip\noindent\mycase{1} $b = \max(E_{odd})$.
	We further consider the following two sub-cases based on the number
	of distinct concentrations in $E$:
	\begin{description}
		\item\mycase{1.1} $\nofconcentrations{E}=3$.
		Then $E=\braced{\nofdroplets{E_{odd}}-2:\min(E_{odd}),\nofdroplets{E_{even}}:a,2:b}$,
		with $\nofdroplets{E_{odd}}\leq 4$.
		If $\nofdroplets{E_{odd}}=3$, then $\min(E_{odd})$ is a singleton, and, 
		after the mixing, $\nofconcentrations{E}=2$ holds,
		which implies that $E$ is now near-final, by Lemma~\ref{lem: mixability m = 2}
		and thus $(b,\min(E_{odd}))$ is actually a near-final pair.		
		Otherwise, $\nofdroplets{E_{odd}}=4$ and $\min(E_{odd})$ is a doubleton.
		This trivially implies that after (and before) the mixing $E$ is near-final,
		so $(b,\min(E_{odd}))$ is a near-final pair.
		(Recall that $|b-\min(E_{odd})|\geq\delta/2$, so the lemma holds
		by Observation~\ref{obs: far-appart mix}.)
				
		
		\item\mycase{1.2} $\nofconcentrations{E}\geq 4$.
		This implies that there is $c\in E_{odd}$ such that $\min(E_{odd})<c<b$;
		mixing $b$ and $c$ produces a non-singleton other than $a$,
		which preserves Invariant~(I), so $(b,c)$ is (I)-safe.
		We analyze the following sub-cases:
		\begin{description}
			\vspace{-0.03in}
			\item\mycase{1.2.1} $c<a$.
			Since $|b-c|\geq\delta/2$, mixing $b$ and $c$ decreases $\potential(E)$ 
			at least by a factor of $1-1/8n$, so we mix them.
			
			\item\mycase{1.2.2} $c>a$.
			Let $d=\half(b+c)$ be the output of the mixing between $b$ and $c$.
			If $d$ is odd, then mixing $d$ and $\min(E_{odd})$ decreases $\potential(E)$ 
			at least by a factor of $1-1/8n$. (Pair $(d, \min(E_{odd}))$ is (I)-safe
			because its mixing produces a non-singleton other than $a$.)
			Otherwise, $d$ is even and, as $|d-a|\geq\delta/4$, 
			mixing $a$ and $d$ decreases $\potential(E)$ at least by a factor of $1-1/32n$.
			(Note that $(a,d)$ is (I)-safe because, after its mixing, $a$ remains non-singleton.)
		\end{description}
		
	\end{description}

	\noindent\mycase{2} $b < \max(E_{odd})$.
	This is similar to Case~$1.2.2$ above under the assumption that $b < \max(E_{odd})$
	and using $c=\max(E_{odd})$.
\end{proof}

%%%%%%%%%%%%%%%%%%%%

\begin{lemma}\label{lem: far mix in E odd non-overlapping E even}
	Assume that $E$ with $\nofdroplets{E}=n\geq 7$ satisfies Invariant~(I).
	Also, assume that $\nofdroplets{E_{\pi}}\geq\nofdroplets{E_{\bar{\pi}}}\geq1$ and
	that $\nofconcentrations{E_\pi} = 1$ with $a\in E_\pi$.
	If either $a < \min(E_{\bar{\pi}})$ or $a>\max(E_{\bar{\pi}})$,	then there exists
	a mixing sequence (of (I)-safe pairs) of length at most $128ns(E) + 1$
	that decreases $\potential(E)$ at least by a factor of $1-1/32\gamma^2n$, 
	for constant $\gamma\in\posintegers$.
\end{lemma}

\begin{proof}
	Let $\delta=\diameter(E)$ before any mixing operation.
	Assume without loss of generality that $\pi=even$ and that $a<\min(E_{odd})$.
	We analyze two sub-cases:
	
	\medskip\noindent\mycase{1} $\diameter(E_{odd})\geq\delta/2\gamma$.
	$E$ satisfying Invariant~(I) implies that there is a non-singleton $b\in E_{odd}$.
	Either $\min(E_{odd})$ or $\max(E_{odd})$ is furthest from $b$, so
	assume without loss of generality that $\min(E_{odd})$ is furthest from $b$.
	As $a<\half(b+\min(E_{odd}))$, pair $(b,\min(E_{odd}))$ is (I)-safe 
	(mixing $b$ and $\min(E_{odd})$	produces a non-singleton other than $a$)
	that satisfies $|b-\min(E_{odd})|\geq\delta/4\gamma$.
	Therefore, mixing $b$ and $\min(E_{odd})$ decreases $\potential(E)$ 
	at least by a factor of $1-1/32\gamma^2n$.
	
	\medskip\noindent\mycase{2} $\diameter(E_{odd})<\delta/2\gamma$.
	This implies that $|\min(E_{odd})-a|\geq\delta/2\gamma$.
	(I)-mix $E_{odd}$ using Lemma~$\ref{lem: E-mix c_pi n >= 22}$;
	$E$ satisfying Invariant~(I) (and our proof in 
	Section~\ref{subsec: proof for n greater-equal than 7}) 
	implies that an even concentration $b$ is eventually produced.
	Since $a$'s frequency is at least $\ceil{n/2}\geq 3$, pair $(a,b)$ is (I)-safe;
	$a$ remains non-singleton after the mixing.
	Additionally, since $b > \min(E_{odd})$, $|b-a|>\delta/2\gamma$ and thus
	mixing $a$ and $b$ decreases $\potential(E)$ at least by a factor of $1-1/8\gamma^2n$.
	Finally, Lemma~$\ref{lem: E-mix c_pi n >= 22}$ takes at most 
	$128ns(E)$ mixing operations, so the lemma holds.
\end{proof}

%%%%%%%%%%%%%%%%%%%%

\begin{lemma}\label{lem: far mix in E even}
	Assume that $E$ with $\nofdroplets{E}=n\geq 22$ satisfies Invariant~$(I)$.
	Also, assume that $\nofdroplets{E_{\pi}}\geq\nofdroplets{E_{\bar{\pi}}}$
	and that $\nofconcentrations{E_{\pi}}\geq 2$.
	If $\diameter(E_\pi) \geq \diameter(E)/\gamma$, for constant $\gamma\in\posintegers$, 
	then there exists an (I)-safe pair whose mixing decreases $\potential(E)$ 
	at least by a factor of $1-1/8\gamma^2n$.
\end{lemma}

\begin{proof}
	Assume without loss of generality that $\pi=even$ and let $a=\min(E_{even})$ and $b=\max(E_{even})$.
	Mixing $a$ and $b$ decreases $\potential(E)$ at least by a factor of $1-1/2\gamma^2n$.
	So, if $(a,b)$ is (I)-safe then we are done.
	Instead, assume otherwise; namely, that $(a,b)$ is not (I)-safe.
	
	The above assumption, and Observation~$\ref{obs: number of same-parity unsafe pairs}$,
	imply that there is $c\in E_{even}$ with $c\notin\braced{a,b}$ 
	for which any pair involving $c$ is (I)-safe.
	Now, either $a$ or $b$ is furthest from $c$, so	assume without loss of generality 
	that $a$ is furthest from $c$.
	Then, $|c-a|\geq\diameter(E)/2\gamma$ and thus mixing $a$ and $c$ decreases 
	$\potential(E)$ at least by a factor of $1-1/8\gamma^2n$.
\end{proof}

%%%%%%%%%%%%%%%%%%%%

\begin{lemma}\label{lem: E-mix c_pi n >= 22}
	Assume that $E$ with $\nofdroplets{E}=n\geq 7$ satisfies Invariant~$(I)$.
	Then, there exists a mixing sequence (of (I)-safe pairs) 
	of length at most $128ns(E)$ that (I)-mixes $E_\pi$.
\end{lemma}

\begin{proof}
	Assume without loss of generality that $\pi = even$ and let $\nofdroplets{E_{even}}=n'$ and
	$\delta=\diameter(E_{even})$ before any mixing operation.
	If a mixing in $E_{even}$ produces droplets with an odd concentration,
	then these droplets are excluded from $E_{even}$ and included in $E_{odd}$;
	this can happen at most $n'/2$ times before $E_{even}$ becomes (I)-mixed.
	
	Assume that the number of (I)-safe pairs in $E_{even}$ is non-zero and 
	let $a=\min(E_{even})$ and $b=\max(E_{even})$.
	If an (I)-safe pair $x,y\in E_{even}$ satisfies $|x-y|\geq \delta/4$,
	then mixing $x$ and $y$ decreases 
	$\potential(E_{even})$ at least by a factor of $1-1/32n'$.
	It follows from Observation~$\ref{obs: far-appart constant mix}$
	that after at most $32n'$ such mixing operations, $\potential(E_{even})$
	decreases at least by a factor of $\half$.
	Hence, $E_{even}$ can be (I)-mixed after at most $32n'\log{\potential(E_{even})}$
	such mixing operations.
	
	We next show that if $E_{even}$ has not been (I)-mixed, then
	after at most two consecutive (I)-safe mixing operations
	either an odd concentration is produced 
	or $\potential(E_{even})$ decreases at least by a factor of $1-1/32n'$.
	Consequently, after at most $128ns(E)$ (I)-safe mixing operations, $E_{even}$ becomes (I)-mixed.
	
	So, if $(a,b)$ is (I)-safe then we are done; $|a-b|= \delta$
	and thus mixing $a$ and $b$ decreases $\potential(E)$ at least by a factor of $1-1/2n$.
	Instead, assume that $(a,b)$ is not (I)-safe and consider the following cases:
	
	\medskip\noindent\mycase{1} $E_{even}$ has only singletons.
	This implies that no mixing in $E_{even}$ decreases the number of non-singletons in $E$.
	So, let $c,d\in E_{even}$ be the furthest-apart (I)-safe pair
	and $x=\half(c+d)$ be the output of their mixing.
	If either $|c-d|\geq \delta/2$ or $x$ odd holds, then we are done.
	Otherwise, let $y\in\braced{a,b}$ be furthest from $x$.
	By the choice of $y$, we have that $|x-y|\geq \delta/2$,
	and as $x$ is a doubleton, $(x,y)$ is (I)-safe and thus we mix it.
	
	\medskip\noindent\mycase{2} $E_{even}$ has exactly one non-singleton $c$.
	Either $a$ or $b$ is furthest from $c$, so assume without loss of generality 
	that $b$ is furthest from $c$; $|b-c|\geq \delta/2$.
	If $(b,c)$ is (I)-safe then we mix it and we are done.
	Otherwise, $(b,c)$ is not (I)-safe, therefore there are exactly two non-singletons in $E$,
	$c\in E_{even}$ (which is a doubleton) and $e=(b+c)/2\notin E_{even}$.
	Since there is an (I)-safe pair in $E_{even}$, there is some $d\in E_{even}$
	such that $d\notin\braced{b,c}$. Then, $(c,d)$ is (I)-safe because 
	mixing $c$ and $d$ produces a non-singleton $x=\half(c+d)$ other than $e$.
	So, mix $c$ and $d$. If $|c-d|\geq\delta/2$ then we are done; similarly if $x$ is odd.
	Otherwise, $(b,x)$ is (I)-safe ($e\neq\half(b+x)$) and $|x-b|\geq\delta/4$, so we mix it.
	
	\medskip\noindent\mycase{3} $E_{even}$ has at least two non-singletons $c<d$.
	We analyze two sub-cases:
	\begin{description}
		\item{\mycase{3.1}} $a=c$.
		Since $(a,b)$ is not (I)-safe, then mixing $b$ and $c$ decreases
		the number of non-singletons down to one. This gives us two sub-cases:
		\begin{description}
			\item{\mycase{3.1.1}} $d=b$.
			$E_{even}$ having an (I)-safe pair implies that there is some $x\in E_{even}$ such that
			$x\notin\braced{c,d}$. Either $c$ or $d$ is furthest from $x$, so			
			assume without loss of generality that $d$ is furthest from $x$; $|d-x|\geq\delta/2$.
			Mixing $x$ and $d$ produces a non-singleton other than $c$, so $(x,d)$ is (I)-safe
			and we mix it.			

			\item{\mycase{3.1.2}} $d<b$.
			This implies that $d=\half(c+b)$, $b-d=\delta/2$ and $(b,d)$ is (I)-safe
			(mixing $b$ and $d$ produces a non-singleton other than $c$), so we mix $b$ and $d$.
		\end{description}

		\item{\mycase{3.2}} $a<c$.
		If $|a-c|\geq\delta/2$ then we mix $a$ and $c$ and we are done;
		$\half(a+c) < d$ implies that mixing $a$ and $c$ produces a non-singleton other than $d$,
		so $(a,c)$ is (I)-safe.
		Instead, let $|a-c|<\delta/2$ which implies $|b-c|\geq\delta/2$.
		Now, if pair $(b,c)$ is (I)-safe then we mix it and we are done, 
		so also assume that $(b,c)$ is not (I)-safe.
		(We cannot have $b=d$ because that would contradict $(a,b)$ not being (I)-safe;
		$|a-c|<\delta/2$ implies $c < \half(a+b)$, so mixing $a$ and $b$ would produce
		a non-singleton other than $c$.)
		Therefore, as $(b,c)$ is not (I)-safe, we have that $d = \half(b+c)$ and mixing $b$ and $d$ 
		produces a non-singleton other than $c$. Hence, $(b,d)$ is (I)-safe
		satisfying $|b-d|\geq\delta/4$, and we mix it.
	\end{description}

	All the pairs mentioned above have difference at least $\delta/4$, so their mixing decreases 
	$\potential(E_{even})$	at least by a factor of $1-1/32n'$.

\end{proof}





We now complete the proof of Theorem~\ref{thm: miscibility characterization}(b).
In Section~\ref{sec: sufficiency of condition mc} we have already established an existence of a perfect-mixing 
sequence for $C$ with intermediate precision $1$.
It still remains to show that this mixing sequence can be modified to have length that is polynomial in $s(C)$.

As explained in the beginning of Section~\ref{sec: polynomial proof}, it remains to
give such a bound for the configuration $E$ constructed from $C$ in Section~\ref{subsec: proof for n greater-equal than 7}
and for mixing sequences with precision $0$ (that is, with only integral concentrations).
Recall that in Section~\ref{sec: polynomial - near final} we gave a polynomial bound if $E$ is near-final.
In Section~\ref{sec: an exponential bound} we gave a bound for arbitrary configurations $E$ that is
exponential only in $n$, thus yielding a polynomial bound for constant $n$.
We will be using these properties in our construction in this section. 
If $E$ is not near-final and $n$ is arbitrary, the general idea for achieving a polynomial bound
is similar to that from Sections~\ref{sec: polynomial - near final} and~\ref{sec: an exponential bound}:
we try to mix a pair of concentrations that are far apart,
to guarantee that we make a quick progress, as measured by the decrease of $\potential(E)$. 
(This is captured by Observation~\ref{obs: far-appart mix}, that will be used frequently and without
explicit reference.) Although this is not posible in each step, we (essentially)
show that for any $E$ there is a polynomial-length sequence of mix operations after which such a
far-apart pair will exist, yielding an overall polynomial bound on a perfect-mixing sequence.
Additionally, to obtain such a perfect-mixing sequence,
at each step we must mix pairs that are either $(\lambda)$-safe, 
for corresponding $\lambda\in\braced{I,I'}$, or near-final.

In this section, for better clarity, we chose to use a top-down presentation style, starting with the statement
of the main theorem (Theorem~\ref{thm: perfect mixability polynomial} below) and its proof, while
all the necessary lemmas that cover different cases from that proof are presented later.
Also, throughout this section, $\gamma$ denotes a small integral constant; for
concreteness we can assume that $\gamma = 2$.

%%%%%%%%%%%%%%%%%

\begin{theorem} \label{thm: perfect mixability polynomial}
	$E$ can be perfectly mixed with precision $0$ 
	in a polynomial number of mixing operations.
\end{theorem} 

\begin{proof}
	If $n = 4$ (and, more generally, if $n$ is a power of $2$), we can simply use 
	Theorem~\ref{them: n power of two} and we are done.
	Thus, we can assume that $n\geq 5$.	
	If $n<22$, we can perfectly mix $E$ using Theorem~\ref{thm: perfect mixability upper bound},
	where the bound on the length of the mixing sequence is polynomial if $n$ is constant.
	Thus in the rest of the argument below we assume that $n \geq 22$.
	
	We show that as long as $E$ satisfies Invariant~(I) (which is the invariant that applies when $n\ge 22$),
	we can find a polynomial-length sequence of mixing operations that will
	decrease $\potential(E)$ by a factor of $1-\Omega(1/n)$. This will be sufficient to prove the theorem.
	
	So assume that $E$ is not yet perfectly mixed and let $\delta=\diameter(E)$. Let $\pi$ be the parity
	for which $\nofdroplets{E_{\pi}}\geq\nofdroplets{E_{\barpi}}$.
	(Recall that we say that $A\subseteq E$ is (I)-mixed if there is no (I)-safe pair in $A$,
	that is, there no pair of distinct concentrations in $A$ whose mixing preserves Invariant~(I) for $E$.
	We will often use this terminology when mixing pairs of concentrations in $E_\pi,E_\barpi\subseteq E$.)
	We consider several cases.
	
	%%%%%%%%%%
	
	\medskip\noindent
	\mycase{1} $\nofconcentrations{E_{\pi}}=1$.
	Let $a\in E_{\pi}$ (that is, $a$ is the only concentration in $E_{\pi}$). 
	Observe that $E_{\barpi}\neq\emptyset$. We have two sub-cases:
	
	\begin{description}

		\item\mycase{1.1} $\min(E_{\bar{\pi}}) < a < \max(E_{\bar{\pi}})$.
		In this case, per Lemma~\ref{lem: farm mix in E odd overlapping E even},
		there are at most two consecutive mixing operations 
		(involving either (I)-safe or near-final pairs)
		after which $\potential(E)$ decreases at least by a factor of $1 - 1/32n$.

		\item \mycase{1.2} $a < \min(E_{\bar{\pi}})$ or $a > \max(E_{\bar{\pi}})$. In this case,
		per Lemma~\ref{lem: far mix in E odd non-overlapping E even}, there is a sequence of
		at most $128ns(E)+1$ (I)-safe mixing operations after which
		$\potential(E)$ decreases at least by a factor of $1 - 1/32\gamma^2n$.
		
	\end{description}	

	%%%%%%%%%%
	
	\medskip\noindent
	\mycase{2} $\nofconcentrations{E_{\pi}}\geq 2$ and $\diameter(E_\pi)\geq\delta/\gamma$.
	In this case, per Lemma~\ref{lem: far mix in E even},	
	there is an (I)-safe pair in $E_{\pi}$ whose mixing decreases
	$\potential(E)$ at least by a factor of $1 - 1/8\gamma^2n$.
	
	%%%%%%%%%%
	
	\medskip\noindent
	\mycase{3} $\nofconcentrations{E_{\pi}}\geq 2$ and $\diameter(E_\pi)<\delta/\gamma$.
	Using Lemma~\ref{lem: E-mix c_pi n >= 22}, there is a sequence of at most
	$128ns(E)$ (I)-safe mixing operations in $E_\pi$, such that after this sequence
	$E_{\pi}$ is (I)-mixed. Let $E'$ be the configuration obtained from $E$
	after applying this sequence. We consider two sub-cases:	
	
	\begin{description}

		\item\mycase{3.1} $\nofdroplets{E'_{\pi}}\geq\nofdroplets{E'_{\bar{\pi}}}$.
		By Observation~\ref{obs: number of same-parity unsafe pairs} below and
		since $E'_{\pi}$ is (I)-mixed, we have that $\nofconcentrations{E'_{\pi}}=1$.		
		Hence, as in Case~$1$ above, there is a sequence of at most $128ns(E')+1$ mixing operations 
		(involving either (I)-safe or near-final pairs)
		after which $\potential(E')$ decreases at least by a factor of $1-1/32\gamma^2n$.
		Thus, $\potential(E)$ also decreases at least by a factor of $1 - 1/32\gamma^2n$.

		\item\mycase{3.2} $\nofdroplets{E'_{\pi}}<\nofdroplets{E'_{\bar{\pi}}}$.
		As $\gamma=2$, $\diameter(E'_\pi)<\delta/\gamma$, and 
		because	the mixing operations on $E_{\pi}$ produced at least one droplet 
		with parity $\bar{\pi}$, we have that 
		$\diameter(E'_{\bar{\pi}})>\delta/\gamma$.
		Hence, as in Case~$2$ above, 
		there is an (I)-safe pair whose mixing decreases
		$\potential(E')$ at least by a factor of $1 - 1/8\gamma^2n$.
		Thus, such mixing also decreases $\potential(E)$ at least by a factor of $1 - 1/8\gamma^2n$.				
	\end{description}
	
	Applying a sequence of mixing operations specified by the cases above
	results in a decrease of $\potential(E)$ at least by a factor of $1 - 1/32\gamma^2n$.
	Thus, by a simple extension to Observation~\ref{obs: far-appart constant mix},
	we obtain that after at most $32\gamma^2n$ such mixing sequences $\potential(E)$ decreases at least by half.
	It follows that after at most $32\gamma^2n\log\potential(E)$ of these mixing sequences (where $\potential(E)$
	denotes the initial potential value), $E$ becomes (I)-mixed, and,
	as $E$ satisfies Invariant~(I), there is a near-final pair in $E$ that we then mix to make $E$ near-final.
	Since $\log\potential(E)\le 2s(E)$, the length of this sequence is at most $64\gamma^2ns(E)$ (where $s(E)$ 
	represents the original size of $E$).
	Each such mixing sequence involves at most $256ns(E)+1$ mixing operations and,
	by~Theorem~\ref{thm: E near-final}, if $E$ is near-final then it can be perfectly mixed 
	by a sequence of at most $144n^3s^2(E)$	mixing operations. 
	Therefore the total number of mixing operations to perfectly mix $E$ is
	at most $2^{14}\gamma^2n^2s^2(E) + 64\gamma^2ns(E) + 144n^3s^2(E)$.
\end{proof}

%%%%%%%%%%%%%%%%%%%%

\begin{observation}\label{obs: number of same-parity unsafe pairs}
	Assume that $E$ with $\nofdroplets{E}=n\geq 22$ satisfies Invariant~(I)
	and let $\pi$ be the parity for which $\nofdroplets{E_{\pi}}\geq\nofdroplets{E_{\bar{\pi}}}$.
	There is at least one droplet in $E_{\pi}$ such that, when paired with any other
	droplet in $E_{\pi}$, the resulting pair is (I)-safe.
\end{observation}

\begin{proof}
	Lemma~$\ref{lem: at most one mod-unsafe pairs}$(a) and the number of distinct odd 
	prime factors of $n$ being less than $\log_3{n}$ imply that $E$ has at most
	$2\lfloor\log_3{n}\rfloor$ droplets that are $\barp$-unsafe when paired with other droplets in $E$.
	This, and Observation~$\ref{obs: decrease number of non-singletons}$ below
	give that the number of droplets that, when mixed with other droplets in $E$
	violate Invariant~(I), is at most 
	$2\lfloor\log_3{n}\rfloor+6< n/2\leq\nofdroplets{E_{\pi}}$, since $n\geq22$.
	Thus, the lemma holds.
\end{proof}
%%%%%%%%%%%%%%%%%%%%

\begin{observation}\label{obs: decrease number of non-singletons}
Assume that $E$ with $\nofdroplets{E}=n\geq 7$ satisfies Invariant~(I).
The number of droplets involved in mixing operations that decrease the number of
non-singletons in $E$ down to one is at most $6$.
\end{observation}

\begin{proof}
	First of all, if the number of non-singletons in $E$ is more than three, 
	then no mixing decreases the number of non-singletons down to one;
	similarly when mixing a non-singleton with frequency higher than two.
	Additionally, mixing two singletons does not decrease the number of non-singletons.	
	Now, $E$ satisfying Invariant~(I) implies that there are 
	at least two non-singletons $a,b\in E$.	
	We consider two types of situations where a mixing decreases the number 
	of non-singletons in $E$ down to one:
	
	\medskip\noindent\mycase{1} Mixing two non-singletons, say $a$ and $b$.
	This can happen when the frequency of both $a$ and $b$ is two each,
	leading to a total of $4$ droplets involved.
	(There could be another non-singleton $e=\half(a+b)$ in $E$, however no mixing 
	involving $e$ would decrease the number of non-singletons down to one
	because either $a$ or $b$ would remain non-singleton after the mixing.)	
	
	\medskip\noindent\mycase{2} Mixing a non-singleton with a singleton.
	(This can happen when $E$ has exactly two non-singletons, $a$ and $b$ respectively.)
	Let $a$ and $c$ be the non-singleton and singleton, respectively.	
	If $a$ is a doubleton and $b=\half(a+c)$, then mixing $a$ and $c$ decreases
	the number of non-singletons down to one.
	Similarly, if $b$ is a doubleton and $a=\half(b+d)$, for some singleton $d\in E$,
	then mixing $b$ and $d$ decreases the number of non-singletons down to one.
	Therefore, the total number of droplets (excluding $a$ and $b$, which were already
	counted in Case~$1$ above) is at most $2$.
	
	\medskip
	Therefore, the number of droplets involved in mixing operations that decrease 
	the number of non-singletons down to one is at most $6$.
\end{proof}

%%%%%%%%%%%%%%%%%%%%

\begin{lemma}\label{lem: farm mix in E odd overlapping E even}
	Assume that $E$ with $\nofdroplets{E}=n\geq 7$ satisfies Invariant~(I).
	Also, assume that $\nofdroplets{E_{\pi}}\geq\nofdroplets{E_{\bar{\pi}}}\geq 1$ and
	that $\nofconcentrations{E_{\pi}} = 1$, with $a\in E_{\pi}$.
	If $\min(E_{\bar{\pi}})<a<\max(E_{\bar{\pi}})$ then there exist at most
	two consecutive mixing operations (involving either (I)-safe or near-final pairs) that
	decrease $\potential(E)$ at least by a factor of $1-1/32n$.
\end{lemma}   

\begin{proof}
	Let $\delta=\diameter(E)$ before any mixing operation.
	Assume without loss of generality that $\pi=even$.
	Since $E$ satisfies Invariant~(I), there is at least one non-singleton $b\in E_{odd}$.
	Either $\min(E_{odd})$ or $\max(E_{odd})$ is furthest from $b$, so
	assume without loss of generality that $\min(E_{odd})$ is furthest from $b$; 
	$|b-\min(E_{odd})|\geq\delta/2$.	
	If pair $(b,\min(E_{odd}))$ is (I)-safe, then we are done.
	So, assume otherwise, namely that pair $(b,\min(E_{odd}))$ is not (I)-safe;
	$a=\half(b+\min(E_{odd}))$ and $b$ (which is a doubleton) are non-singletons in $E$.
	(There could be at most one more non-singleton in $E$, namely $\min(E_{odd})$ strictly doubleton;
	otherwise $(b,\min(E_{odd}))$ would be in fact (I)-safe.)
	Consider the following cases:
	
	\medskip\noindent\mycase{1} $b = \max(E_{odd})$.
	We further consider the following two sub-cases based on the number
	of distinct concentrations in $E$:
	\begin{description}
		\item\mycase{1.1} $\nofconcentrations{E}=3$.
		Then $E=\braced{\nofdroplets{E_{odd}}-2:\min(E_{odd}),\nofdroplets{E_{even}}:a,2:b}$,
		with $\nofdroplets{E_{odd}}\leq 4$.
		If $\nofdroplets{E_{odd}}=3$, then $\min(E_{odd})$ is a singleton, and, 
		after the mixing, $\nofconcentrations{E}=2$ holds,
		which implies that $E$ is now near-final, by Lemma~\ref{lem: mixability m = 2}
		and thus $(b,\min(E_{odd}))$ is actually a near-final pair.		
		Otherwise, $\nofdroplets{E_{odd}}=4$ and $\min(E_{odd})$ is a doubleton.
		This trivially implies that after (and before) the mixing $E$ is near-final,
		so $(b,\min(E_{odd}))$ is a near-final pair.
		(Recall that $|b-\min(E_{odd})|\geq\delta/2$, so the lemma holds
		by Observation~\ref{obs: far-appart mix}.)
				
		
		\item\mycase{1.2} $\nofconcentrations{E}\geq 4$.
		This implies that there is $c\in E_{odd}$ such that $\min(E_{odd})<c<b$;
		mixing $b$ and $c$ produces a non-singleton other than $a$,
		which preserves Invariant~(I), so $(b,c)$ is (I)-safe.
		We analyze the following sub-cases:
		\begin{description}
			\vspace{-0.03in}
			\item\mycase{1.2.1} $c<a$.
			Since $|b-c|\geq\delta/2$, mixing $b$ and $c$ decreases $\potential(E)$ 
			at least by a factor of $1-1/8n$, so we mix them.
			
			\item\mycase{1.2.2} $c>a$.
			Let $d=\half(b+c)$ be the output of the mixing between $b$ and $c$.
			If $d$ is odd, then mixing $d$ and $\min(E_{odd})$ decreases $\potential(E)$ 
			at least by a factor of $1-1/8n$. (Pair $(d, \min(E_{odd}))$ is (I)-safe
			because its mixing produces a non-singleton other than $a$.)
			Otherwise, $d$ is even and, as $|d-a|\geq\delta/4$, 
			mixing $a$ and $d$ decreases $\potential(E)$ at least by a factor of $1-1/32n$.
			(Note that $(a,d)$ is (I)-safe because, after its mixing, $a$ remains non-singleton.)
		\end{description}
		
	\end{description}

	\noindent\mycase{2} $b < \max(E_{odd})$.
	This is similar to Case~$1.2.2$ above under the assumption that $b < \max(E_{odd})$
	and using $c=\max(E_{odd})$.
\end{proof}

%%%%%%%%%%%%%%%%%%%%

\begin{lemma}\label{lem: far mix in E odd non-overlapping E even}
	Assume that $E$ with $\nofdroplets{E}=n\geq 7$ satisfies Invariant~(I).
	Also, assume that $\nofdroplets{E_{\pi}}\geq\nofdroplets{E_{\bar{\pi}}}\geq1$ and
	that $\nofconcentrations{E_\pi} = 1$ with $a\in E_\pi$.
	If either $a < \min(E_{\bar{\pi}})$ or $a>\max(E_{\bar{\pi}})$,	then there exists
	a mixing sequence (of (I)-safe pairs) of length at most $128ns(E) + 1$
	that decreases $\potential(E)$ at least by a factor of $1-1/32\gamma^2n$, 
	for constant $\gamma\in\posintegers$.
\end{lemma}

\begin{proof}
	Let $\delta=\diameter(E)$ before any mixing operation.
	Assume without loss of generality that $\pi=even$ and that $a<\min(E_{odd})$.
	We analyze two sub-cases:
	
	\medskip\noindent\mycase{1} $\diameter(E_{odd})\geq\delta/2\gamma$.
	$E$ satisfying Invariant~(I) implies that there is a non-singleton $b\in E_{odd}$.
	Either $\min(E_{odd})$ or $\max(E_{odd})$ is furthest from $b$, so
	assume without loss of generality that $\min(E_{odd})$ is furthest from $b$.
	As $a<\half(b+\min(E_{odd}))$, pair $(b,\min(E_{odd}))$ is (I)-safe 
	(mixing $b$ and $\min(E_{odd})$	produces a non-singleton other than $a$)
	that satisfies $|b-\min(E_{odd})|\geq\delta/4\gamma$.
	Therefore, mixing $b$ and $\min(E_{odd})$ decreases $\potential(E)$ 
	at least by a factor of $1-1/32\gamma^2n$.
	
	\medskip\noindent\mycase{2} $\diameter(E_{odd})<\delta/2\gamma$.
	This implies that $|\min(E_{odd})-a|\geq\delta/2\gamma$.
	(I)-mix $E_{odd}$ using Lemma~$\ref{lem: E-mix c_pi n >= 22}$;
	$E$ satisfying Invariant~(I) (and our proof in 
	Section~\ref{subsec: proof for n greater-equal than 7}) 
	implies that an even concentration $b$ is eventually produced.
	Since $a$'s frequency is at least $\ceil{n/2}\geq 3$, pair $(a,b)$ is (I)-safe;
	$a$ remains non-singleton after the mixing.
	Additionally, since $b > \min(E_{odd})$, $|b-a|>\delta/2\gamma$ and thus
	mixing $a$ and $b$ decreases $\potential(E)$ at least by a factor of $1-1/8\gamma^2n$.
	Finally, Lemma~$\ref{lem: E-mix c_pi n >= 22}$ takes at most 
	$128ns(E)$ mixing operations, so the lemma holds.
\end{proof}

%%%%%%%%%%%%%%%%%%%%

\begin{lemma}\label{lem: far mix in E even}
	Assume that $E$ with $\nofdroplets{E}=n\geq 22$ satisfies Invariant~$(I)$.
	Also, assume that $\nofdroplets{E_{\pi}}\geq\nofdroplets{E_{\bar{\pi}}}$
	and that $\nofconcentrations{E_{\pi}}\geq 2$.
	If $\diameter(E_\pi) \geq \diameter(E)/\gamma$, for constant $\gamma\in\posintegers$, 
	then there exists an (I)-safe pair whose mixing decreases $\potential(E)$ 
	at least by a factor of $1-1/8\gamma^2n$.
\end{lemma}

\begin{proof}
	Assume without loss of generality that $\pi=even$ and let $a=\min(E_{even})$ and $b=\max(E_{even})$.
	Mixing $a$ and $b$ decreases $\potential(E)$ at least by a factor of $1-1/2\gamma^2n$.
	So, if $(a,b)$ is (I)-safe then we are done.
	Instead, assume otherwise; namely, that $(a,b)$ is not (I)-safe.
	
	The above assumption, and Observation~$\ref{obs: number of same-parity unsafe pairs}$,
	imply that there is $c\in E_{even}$ with $c\notin\braced{a,b}$ 
	for which any pair involving $c$ is (I)-safe.
	Now, either $a$ or $b$ is furthest from $c$, so	assume without loss of generality 
	that $a$ is furthest from $c$.
	Then, $|c-a|\geq\diameter(E)/2\gamma$ and thus mixing $a$ and $c$ decreases 
	$\potential(E)$ at least by a factor of $1-1/8\gamma^2n$.
\end{proof}

%%%%%%%%%%%%%%%%%%%%

\begin{lemma}\label{lem: E-mix c_pi n >= 22}
	Assume that $E$ with $\nofdroplets{E}=n\geq 7$ satisfies Invariant~$(I)$.
	Then, there exists a mixing sequence (of (I)-safe pairs) 
	of length at most $128ns(E)$ that (I)-mixes $E_\pi$.
\end{lemma}

\begin{proof}
	Assume without loss of generality that $\pi = even$ and let $\nofdroplets{E_{even}}=n'$ and
	$\delta=\diameter(E_{even})$ before any mixing operation.
	If a mixing in $E_{even}$ produces droplets with an odd concentration,
	then these droplets are excluded from $E_{even}$ and included in $E_{odd}$;
	this can happen at most $n'/2$ times before $E_{even}$ becomes (I)-mixed.
	
	Assume that the number of (I)-safe pairs in $E_{even}$ is non-zero and 
	let $a=\min(E_{even})$ and $b=\max(E_{even})$.
	If an (I)-safe pair $x,y\in E_{even}$ satisfies $|x-y|\geq \delta/4$,
	then mixing $x$ and $y$ decreases 
	$\potential(E_{even})$ at least by a factor of $1-1/32n'$.
	It follows from Observation~$\ref{obs: far-appart constant mix}$
	that after at most $32n'$ such mixing operations, $\potential(E_{even})$
	decreases at least by a factor of $\half$.
	Hence, $E_{even}$ can be (I)-mixed after at most $32n'\log{\potential(E_{even})}$
	such mixing operations.
	
	We next show that if $E_{even}$ has not been (I)-mixed, then
	after at most two consecutive (I)-safe mixing operations
	either an odd concentration is produced 
	or $\potential(E_{even})$ decreases at least by a factor of $1-1/32n'$.
	Consequently, after at most $128ns(E)$ (I)-safe mixing operations, $E_{even}$ becomes (I)-mixed.
	
	So, if $(a,b)$ is (I)-safe then we are done; $|a-b|= \delta$
	and thus mixing $a$ and $b$ decreases $\potential(E)$ at least by a factor of $1-1/2n$.
	Instead, assume that $(a,b)$ is not (I)-safe and consider the following cases:
	
	\medskip\noindent\mycase{1} $E_{even}$ has only singletons.
	This implies that no mixing in $E_{even}$ decreases the number of non-singletons in $E$.
	So, let $c,d\in E_{even}$ be the furthest-apart (I)-safe pair
	and $x=\half(c+d)$ be the output of their mixing.
	If either $|c-d|\geq \delta/2$ or $x$ odd holds, then we are done.
	Otherwise, let $y\in\braced{a,b}$ be furthest from $x$.
	By the choice of $y$, we have that $|x-y|\geq \delta/2$,
	and as $x$ is a doubleton, $(x,y)$ is (I)-safe and thus we mix it.
	
	\medskip\noindent\mycase{2} $E_{even}$ has exactly one non-singleton $c$.
	Either $a$ or $b$ is furthest from $c$, so assume without loss of generality 
	that $b$ is furthest from $c$; $|b-c|\geq \delta/2$.
	If $(b,c)$ is (I)-safe then we mix it and we are done.
	Otherwise, $(b,c)$ is not (I)-safe, therefore there are exactly two non-singletons in $E$,
	$c\in E_{even}$ (which is a doubleton) and $e=(b+c)/2\notin E_{even}$.
	Since there is an (I)-safe pair in $E_{even}$, there is some $d\in E_{even}$
	such that $d\notin\braced{b,c}$. Then, $(c,d)$ is (I)-safe because 
	mixing $c$ and $d$ produces a non-singleton $x=\half(c+d)$ other than $e$.
	So, mix $c$ and $d$. If $|c-d|\geq\delta/2$ then we are done; similarly if $x$ is odd.
	Otherwise, $(b,x)$ is (I)-safe ($e\neq\half(b+x)$) and $|x-b|\geq\delta/4$, so we mix it.
	
	\medskip\noindent\mycase{3} $E_{even}$ has at least two non-singletons $c<d$.
	We analyze two sub-cases:
	\begin{description}
		\item{\mycase{3.1}} $a=c$.
		Since $(a,b)$ is not (I)-safe, then mixing $b$ and $c$ decreases
		the number of non-singletons down to one. This gives us two sub-cases:
		\begin{description}
			\item{\mycase{3.1.1}} $d=b$.
			$E_{even}$ having an (I)-safe pair implies that there is some $x\in E_{even}$ such that
			$x\notin\braced{c,d}$. Either $c$ or $d$ is furthest from $x$, so			
			assume without loss of generality that $d$ is furthest from $x$; $|d-x|\geq\delta/2$.
			Mixing $x$ and $d$ produces a non-singleton other than $c$, so $(x,d)$ is (I)-safe
			and we mix it.			

			\item{\mycase{3.1.2}} $d<b$.
			This implies that $d=\half(c+b)$, $b-d=\delta/2$ and $(b,d)$ is (I)-safe
			(mixing $b$ and $d$ produces a non-singleton other than $c$), so we mix $b$ and $d$.
		\end{description}

		\item{\mycase{3.2}} $a<c$.
		If $|a-c|\geq\delta/2$ then we mix $a$ and $c$ and we are done;
		$\half(a+c) < d$ implies that mixing $a$ and $c$ produces a non-singleton other than $d$,
		so $(a,c)$ is (I)-safe.
		Instead, let $|a-c|<\delta/2$ which implies $|b-c|\geq\delta/2$.
		Now, if pair $(b,c)$ is (I)-safe then we mix it and we are done, 
		so also assume that $(b,c)$ is not (I)-safe.
		(We cannot have $b=d$ because that would contradict $(a,b)$ not being (I)-safe;
		$|a-c|<\delta/2$ implies $c < \half(a+b)$, so mixing $a$ and $b$ would produce
		a non-singleton other than $c$.)
		Therefore, as $(b,c)$ is not (I)-safe, we have that $d = \half(b+c)$ and mixing $b$ and $d$ 
		produces a non-singleton other than $c$. Hence, $(b,d)$ is (I)-safe
		satisfying $|b-d|\geq\delta/4$, and we mix it.
	\end{description}

	All the pairs mentioned above have difference at least $\delta/4$, so their mixing decreases 
	$\potential(E_{even})$	at least by a factor of $1-1/32n'$.

\end{proof}



